\documentclass[11pt,a4paper]{article}
\usepackage[margin=24mm,headheight=15pt]{geometry}
\usepackage[T1]{fontenc}
\usepackage[utf8]{inputenc}
\usepackage{lmodern,microtype}
\usepackage{amsmath,amssymb,amsthm,mathtools}
\usepackage{booktabs,array,longtable,enumitem}
\usepackage{xcolor,fancyhdr}
\usepackage[colorlinks=true,linkcolor=blue!45!black,urlcolor=blue!45!black,citecolor=blue!45!black]{hyperref}
\newtheorem{theorem}{Theorem}
\newtheorem{lemma}[theorem]{Lemma}
\newtheorem{proposition}[theorem]{Proposition}
\newtheorem{corollary}[theorem]{Corollary}
\theoremstyle{remark}
\newtheorem{remark}[theorem]{Remark}
\newcommand{\formalname}[1]{\leavevmode\par\noindent\textnormal{\small Proposed name: \texttt{#1}.}\par\smallskip}
\newcommand{\WG}{\operatorname{wordGrid}}
\newcommand{\can}{\operatorname{can}}
\setlist{itemsep=3pt,topsep=5pt}
\setlength{\emergencystretch}{3em}
\setlength{\parskip}{3pt}
\pagestyle{fancy}
\fancyhf{}
\fancyhead[L]{\small Statement E: the exterior native map}
\fancyfoot[C]{\thepage}
\title{Statement E: An Explicit Exterior Native Map\\[6pt]
\large A finite event construction and a direct canonical-class proof}
\author{}
\date{2 November 2026}
\begin{document}
\maketitle
\begin{abstract}
We prove Statement E of the supplied problem brief, using its expressly
granted analytic word, native-complex, and binary-seam inputs.  A finite
projection of a suitably chosen nearby normalization gives a transitive
factorization of the central event word into conjugated local branch
words and positive single generators.  Explicit signed word-grid
relations, a finite Hurwitz algorithm, and star scissors turn this into
an exact seam of marked factor grids.  We give the protected-summand
cancellation rules needed to remove every conjugator.

The local maps are constructed directly on the native complexes over
$\mathbb F_2[T,T^{-1}]$.  A rectangle identity gives a polynomial
degree-zero detecting equivalence for each degree-minus-one O switch.
Component derivatives turn the global class calculation into the
determinant of an explicitly chosen rooted-tree incidence matrix.
Explicit auxiliary contractions and an explicit simple-factor
functional finish the construction.  The resulting map has Maslov
degree zero, shifts $(0,\mu(V)/2)$, and sends the canonical class to
the canonical class.  All additional analytic and combinatorial
lemmas are proved, including the native maps used on intermediate
auxiliary marking pairs.  The document includes exact seam formulas,
the full operation count, exceptional cases, a fully indexed
$S=4,\mu(V)=1$ example, and a dependency ledger.
\end{abstract}
\begingroup
\small
\setlength{\parskip}{0pt}
\tableofcontents
\endgroup
\newpage


\section{Statement E and the inputs used}
\label{sec:statement}

We use the notation and the explicitly granted inputs of the supplied
problem brief~\cite{brief}.  An \emph{assembly} and a \emph{level} mean
the data specified there, including their invariant identities.  The
assertion proved here is the chain-map statement itself.

\begin{theorem}[Statement E]\label{thm:E}
Let $\boldsymbol n:(\mathbb C^2,0)\to(\mathbb C^3,0)$ be a holomorphic
map germ with an injective representative and $(d\boldsymbol n)_0=0$.
There exist an assembly, an allowed nonzero level $s_0$, a positive word
$\mathbf w$ on $S$ strands, positive words $\beta_j$ on $a'_j$ strands
for the enumerated singular points $Q_j$, and grids $D_j$ such that
\begin{align*}
 &\operatorname{cycles}(\pi(\mathbf w))=b,
 &&\ell(\mathbf w)+S(S-1)+b=2\delta(W)+S,\\
 &2\le a'_j\le t_J,\qquad \operatorname{cycles}(\pi(\beta_j))=1,
 &&\ell(\beta_j)+1=2\delta_j+a'_j,
\end{align*}
each $\beta_j$ contains every generator, and
$\operatorname{wordGrid}(\beta_j)\sim D_j$.
Every $D_j$ has size at least two.  They are seam prepared, and there
is an $\mathbb F_2$-linear chain map in each Maslov degree
\[
 \psi:C\bigl(\operatorname{wordGrid}(\mathbf d_S\mathbf w)\bigr)
       \longrightarrow C\bigl(\operatorname{iteratedSeamGrid}(D)\bigr)
\]
of Maslov degree zero, with bifiltration shifts $(0,\mu(V)/2)$, for which
\[
 \psi_*\operatorname{can}
       \bigl(\operatorname{wordGrid}(\mathbf d_S\mathbf w)\bigr)
   =\operatorname{can}\bigl(\operatorname{iteratedSeamGrid}(D)\bigr).
\]
One can take $a'_j=a_j$.  More strongly, the constructed $D_j$ satisfy
$X(0)=0$, $O(\mathrm{last})=0$, and
$X(\mathrm{last})=\mathrm{last}$ simultaneously.  Thus all row
rotations in the final iterated seam are the identity.
\end{theorem}

The proof constructs the map directly.  Some intermediate homogeneous
maps have degrees $-1$ and $+1$; their composite has degree zero.
Auxiliary pairs of marking permutations are used for the formulas
of those maps.  The source and target in the theorem are valid native
grids.  This is an alternative to the restricted M/X/B movie in
Section~3.4 of the brief.  In particular, it makes no claim to be a
movie consisting only of those three permitted kinds.

Here is the numerical structure of the construction.  A finite event
factorization has $N$ singular factors and $B$ positive single-generator
factors.  Put
\begin{equation}\label{eq:global-counts}
 r=N+B,\qquad s=(S-1)(r-1),\qquad
 U=\sum_{j=1}^N(S-a_j)+B(S-2).
\end{equation}
The analytic calculation proves
\begin{equation}\label{eq:B-intro}
 B=S-\sum_j a_j+N-1+\mu(V),
 \qquad s-U=\mu(V).
\end{equation}
The scissor construction has $s-U$ X switches and $U$ O switches.
The O switches are paired algebraically with $U$ auxiliary
contractions.  A determinant calculation, proved below from rectangle
counts, gives the canonical-class identity for this complete
composition.  A separate degree-zero normalization of each four-mark
auxiliary block contributes no filtration shift.  Finally an explicit
functional removes each single-generator tensor factor.

\subsection{The admitted inputs}

We use the following statements precisely in the form granted in
Sections~2 and~5 of the brief:
\begin{enumerate}
\item The assembly and arbitrarily small allowed levels exist;
$v=z_0\circ\boldsymbol n$ is reduced; the nearby normalization is
smooth; its singular image germs are unibranch with
$2\le a_j\le t_J$; $N\ge1$; and
\[
 2\delta(V)+M=S,\quad
 \delta(W)=\delta(V)+\sum_j\delta_j,\quad
 \mu(V)=2\delta(V)-b+1,\quad S\ge4.
\]
The branch count of the source and central image is the same $b$.
\item The central and local \emph{event-word} constructions (I1),
(I2), (B1), and (B2) are available for transverse projections.
Thus the central word is $\mathbf d_S\mathbf w$, with
$\mathbf w$ positive of cycle count $b$, and
\[
 \ell(\mathbf d_S\mathbf w)=2\delta(W)+S-b,\qquad
 \ell(\beta_j)=2\delta_j+a_j-1.
\]
The local $\beta_j$ is a positive word with one cycle and every
generator.  We use these event words, rather than an inference from
ambient link isotopy.
\item The native rectangle differential squares to zero.  The
degree-zero polynomial homology image of every valid grid is one
dimensional and nonzero.  The given native commutation and cyclic
permutation maps have degree zero, shifts $(0,0)$, and preserve its
nonzero element.  The supplied combinatorial corner reduction
expresses the other stabilization orientations using the primitive
orientations, commutations, and cyclic permutations.
\item The binary seam--tensor comparison is a filtered chain homotopy
equivalence of degree zero and shifts $(0,0)$ in both directions,
with the stipulated canonical-class identity.  No naturality of
that comparison for derivatives of O variables is assumed or used.
\end{enumerate}
All projection, event-factorization, signed-word, marked-cancellation,
seam-coordinate, auxiliary-map, and class-pairing assertions beyond
these inputs are proved below.  Weierstrass preparation and division,
continuity of polynomial roots, the implicit function theorem,
Puiseux parametrization, and the one-variable maximum principle are
the elementary analytic tools permitted by the brief.  The explicit
compact vector-field flows used below also give their own local
trivializations.

\section{Conventions and native maps along ordinary moves}
\label{sec:conventions}

All coefficients are in $\mathbb F_2$ and
$\Lambda=\mathbb F_2[T,T^{-1}]$.  A size-$n$ grid is specified by
permutations $O,X$ of $\{0,\ldots,n-1\}$, with
$O(c)\ne X(c)$ and $n\ge2$.  A marking in square $(c,r)$ has centre
$(c+\tfrac12,r+\tfrac12)$.  A state $x$ is a permutation, represented
by the lattice points $(c,x(c))$.  For finite sets,
\[
 I(P,Q)=\#\{(p,q):p_1<q_1,\ p_2<q_2\},\qquad
 J(P,Q)=\frac{I(P,Q)+I(Q,P)}2.
\]
Write
\[
 M_O(x)=J(x,x)-2J(x,O)+J(O,O)+1,\qquad
 A(x)=\frac{M_O(x)-M_X(x)}2-\frac{n-1}2 .
\]
The complex is the free $\Lambda$-module on the states, with
\[
 d x=\sum_y\ \sum_{q\in\operatorname{Rect}^{\circ}(x,y)}
            T^{n_O(q)}y,\qquad
 \deg(T^kx)=M_O(x)-2k.
\]
Here the southwest and northeast corners of a rectangle from $x$
to $y$ belong to $x$, its other corners belong to $y$, and its
interior contains no point of $x$.  Rectangles containing X markings
are counted.  The two filtration coordinates of $T^kx$ are
$(-k,A(x)-k)$.  The polynomial subcomplex $P$ is spanned by the
monomials with $k\ge0$.  Its nonzero degree-zero homology image in
$C$ is denoted $\operatorname{can}$.

The following two rectangle formulas follow by the four-corner
indicator calculation in $J$:
\begin{equation}\label{eq:rectangle-gradings}
 M_O(x)-M_O(y)=1-2n_O(q),\qquad
 A(x)-A(y)=n_X(q)-n_O(q)
\end{equation}
for an empty rectangle $q:x\to y$.  For a nonempty rectangle, the
first right-hand side has the additional term
$2\#(\operatorname{Int}q\cap x)$.  To verify the formulas, replace
the two diagonal state points by the other diagonal.  A state point
inside contributes two to the $J(x,x)$ difference; a marking inside
contributes two to the $2J(x,O)$ difference; points outside contribute
zero.  The moving corners contribute one.  Choose cuts outside a
torus rectangle to obtain the same calculation for wrapping
rectangles.  Transposing both coordinates preserves these formulas
and the gradings, with the state permutation replaced by its inverse.

An \emph{auxiliary pair} permits $O(c)=X(c)$.  Its complex and
gradings use the same formulas.  Its canonical class is well defined
whenever $n\ge2$: choose the disjoint permutation $Y(c)=O(c+1)$ and
identify its O complex with the valid grid $(O,Y)$.  The complex,
degree, polynomial subcomplex, and resulting class do not depend
on this choice.  Auxiliary pairs are used only for explicitly given
algebraic maps; the theorem's endpoints are valid grids.

\begin{lemma}[Polynomial nonvanishing; uses the canonical-class input]
\label{lem:N}
\formalname{polynomialImageNonzero}
A degree-zero chain map with first shift zero sends the canonical
class either to zero or to the target's canonical class.  Thus
nonvanishing of that image suffices for the required class identity.
\end{lemma}
\begin{proof}
The map takes $P$ into $P$ and therefore takes the degree-zero
polynomial homology image into that of the target.  The latter has
exactly one nonzero element.  The same argument applies to auxiliary
pairs by their O-complex identification.
\end{proof}

\begin{lemma}[Transport along ordinary moves; uses the native move inputs]
\label{lem:move-transport}
\formalname{movesToNativeMap}
Every explicit finite ordinary-move sequence between grids of size
at least two gives $\Lambda$-linear maps in both directions of degree zero, shifts
$(0,0)$, and canonical-class identity.  After O markings are coloured
constantly on each component, the maps may be chosen to intertwine,
on homology, the sums of the O derivatives of every colour.  A
stabilization gives its new O the colour of its component.
\end{lemma}
\begin{proof}
For the degree, filtration, and class assertions, apply the granted
maps to each commutation and cyclic permutation.  For a primitive
stabilization use the explicit maps verified in
Section~\ref{sec:type-r-independent}; for another corner apply the
supplied combinatorial corner reduction and compose those maps.
Reverse the reduction to obtain the other direction.  Induction on
the length of the sequence proves the assertions.  The separate-variable
proof in Section~\ref{sec:type-r-independent} verifies the derivative
claim for the primitive maps, commutations, and permutations.
Its differentiated identities compose, so the same induction proves
that assertion.  This proves the map version of the corner reduction
and the finite-sequence assertion M702; no invariance theorem for
isotopic links is used.
\end{proof}

For completeness, the seam of sizes $n+1,m+1$, under the conditions
$X_G(n)=n$, $X_H(0)=0$, $O_H(m)=0$, is the size-$n+m+1$ pair
\begin{equation}\label{eq:seam-definition}
 \Sigma.Y(c)=
 \begin{cases}
   G.Y(c),&c<n,\\
   n+H.Y(d),&c=n+d,\ H.Y(d)\ne0,\\
   G.Y(n),&c=n+d,\ H.Y(d)=0,
 \end{cases}
 \quad Y\in\{O,X\}.
\end{equation}
For arbitrary prepared factors, at an intermediate size $q$ the
specified row rotation is addition of
$q-1-X(\mathrm{last})$ modulo $q$.  The iterated size is
$\sum_j\operatorname{size}(D_j)-(N-1)$.  Our stronger preparation
makes each of these additions zero.

Generators have indices starting at zero and
$\rho_m=\sigma_0\cdots\sigma_{m-2}$,
$\mathbf d_m=\rho_m^m$.  The positive letter $\sigma_r$ has the
three-column marking heights
\[
 O=(r+\tfrac12,r,r+1),\qquad
 X=(r,r+1,r+\tfrac12).
\]
The negative letter uses
$O=(r+\tfrac12,r+1,r)$,
$X=(r+1,r,r+\tfrac12)$.
Integer-height arcs join successive occurrences of that rank
cyclically; the middle arc joins the first and third column of its
letter.  Rows are ordered by their heights, with a specified tie
order.  This recovers the positive word grid of the brief up to its
allowed tie changes.  A completely unused rank, in the signed
extension, is retained as an explicit four-mark $2\times2$ component
when it arises by cancellation.  All initial word normalizations
will retain every generator.

Tracing a component through either three-column gadget exchanges
its two integer ranks.  Consequently the components of a closed
signed word grid correspond to the cycles of its word permutation.
This can also be checked by iterating
$c\mapsto X^{-1}(O(c))$: the internal middle arc is traversed once
and the next integer arc is precisely the outgoing arc of the
transposed rank.  This elementary incidence observation will be used
to count the auxiliary components.


\section{A finite projection and its explicit event factorization}
\label{sec:finite-projection-complete}

\begin{lemma}[Finite event factorization; new analytic]
\label{lem:analytic-factorization}
\formalname{finiteProjectionEventFactorization}
For an assembly as in Theorem~\ref{thm:E}, there are an allowed level
$s\ne0$ and a transverse projection $h$ such that a connected compact
normalization level has a proper finite projection of degree $S$,
unramified at its boundary and containing all the level's singular
points.  There is a finite signed-word derivation, up to whole-word
conjugacy, of
\[
 \mathbf d_S\mathbf w\ \sim\
 \prod_{l=1}^{N+B}\gamma_l\eta_l\gamma_l^{-1},
\]
where the cores consist of the $N$ supplied local positive event
words and $B$ positive single generators.  The permutations of the
conjugated factors $\gamma_l\eta_l\gamma_l^{-1}$
generate a transitive group, and
$B=S-\sum_j a_j+N-1+\mu(V)$.  All projection, level, path, and
finite-event choices are made by the rules in the proof.
\end{lemma}
\begin{proof}

We prove the analytic assertions needed for the event construction.  In
particular we do not require that a single linear projection have only
simple critical points or distinct critical values.  Ordinary higher
ramification will be replaced by a finite local family of simple
ramifications, at the level of boundary configurations.

\subsection{An elementary local finite-map fact}

We first spell out the finite-map consequence of Weierstrass preparation
that will be used.  Suppose $F=(f,g)$ is holomorphic near a point of
$\mathbb C^2$ and that its fibre through that point is isolated there.
Then $F$ is open at that point and has finite fibres on a sufficiently
small representative.  Translate the point and its value to zero, and
choose source coordinates $(u,w)$ in which $f$ is regular in $w$.
Preparation, with $s$ as parameter, writes $f(u,w)-s$ as a unit times a
monic polynomial $P(u,s;w)$ of some degree $p$.  Divide $g(u,w)-t$ by
$P$, obtaining a polynomial remainder $Q(u,s,t;w)$ of degree less than
$p$.  Multiplication by $Q$ in the free module with basis
$1,w,\ldots,w^{p-1}$ modulo $P$ is a $p\times p$ matrix.  Its determinant
$R(u,s,t)$ vanishes exactly when $P$ and $Q$ have a common root.  This
last assertion follows by factoring $P$ over $\mathbb C$: multiplication
by $Q$ is invertible in each factor $\mathbb C[w]/(w-a)^e$ exactly when
$Q(a)\ne0$.

The isolated-fibre assumption says that $R(u,0,0)$ has an isolated zero
at $u=0$.  A second preparation therefore writes $R$ as a unit times a
monic polynomial in $u$, with coefficients holomorphic in $(s,t)$.
For every sufficiently small $(s,t)$, its roots remain near zero by
continuity of roots.  Each is accompanied by a common root of $P,Q$,
also near zero.  Thus every sufficiently small target point has a
preimage in any prescribed sufficiently small source neighbourhood.
The same two monic equations bound the number of points in a fibre.
This proves the asserted local openness and finiteness, using only
preparation, division, and continuity of polynomial roots.

We will also use the elementary observation that a compact analytic
subset of an open subset of $\mathbb C^2$ is finite.  If it had a
positive-dimensional irreducible component, choose a coordinate linear
function nonconstant on that component and a point where its absolute
value is maximal.  A local curve branch through that point, parametrized
by preparation or Puiseux expansion, contradicts the one-variable
maximum principle.  A function constant on a local branch of an
irreducible component is constant on that component by analytic
continuation, so the chosen nonconstant function supplies such a
nonconstant branch.  A two-dimensional component contains an open set,
where the usual maximum principle gives the same contradiction.  Finally,
a compact zero-dimensional analytic set is finite, since its points
are isolated and a finite collection of their isolating neighbourhoods
covers it.

\subsection{A proper representative with the correct number of sheets}

Put $v=z_0\circ\boldsymbol n$.  Choose a target-section coordinate
$\xi_0$ transverse to all the central branches, and a complementary
coordinate $\zeta$.  Write $h_0=\xi_0\circ\boldsymbol n$ and
$k=\zeta\circ\boldsymbol n$.  The common zero of $(v,h_0)$ is isolated:
on $v=0$ its image lies in $W\cap\{\xi_0=0\}$, whose germ is just the
origin by transversality; injectivity then gives the unique preimage.

Choose a closed Euclidean source ball $\overline B_R$ inside a
representative on which $\boldsymbol n$ is injective, with zero as the
only common zero of $(v,h_0)$.  Choose it also inside a representative
of the finitely many curve branches and isolated points of the
critical-set germ of $\boldsymbol n$; the finite-branch assertion is
proved below.  We also take it inside the isolated-critical-point
representative of $v$, justified below.  Thus those
later arguments will not require changing
this already fixed ball.  Then
\[
 \eta=\min_{\partial B_R}\max(|v|,|h_0|)>0.
\]
There is no additional boundary-choice assertion hidden here: the
boundary is compact and contains no common zero.  Choose positive
target radii smaller than $\eta$.  The set
\[
 \Omega=\{p\in B_R:|v(p)|<\epsilon_0, |h_0(p)|<\rho_0\}
\]
maps properly to $\Delta_{\epsilon_0}\times\Delta_{\rho_0}$, since the
preimage of any compact subset is a closed subset of $\overline B_R$
disjoint from its boundary.  Its fibres are compact analytic subsets
and hence finite.  Local openness follows from the preceding
Weierstrass argument.

Each connected component of $\Omega$ is open and closed in $\Omega$.
Its image is open by local openness and closed by properness.  Thus
every nonempty component maps onto the connected target bidisc and
contains a point above $(0,0)$.  Since that fibre consists only of zero,
$\Omega$ is connected and the map is onto.

For a small nonzero $\xi_0$, the central curve has exactly $S$ distinct
points above it.  Indeed its transverse Weierstrass degree is its
multiplicity $S$, and on each normalized branch
$\xi_0=t^m u(t)$ with $u(0)\ne0$; after replacing $t$ by
$t\,u(t)^{1/m}$, this is $t^m$.  Its derivative is nonzero off zero,
and the distinct central branches do not intersect away from zero after
shrinking.  Injectivity identifies these $S$ points with their $S$
source preimages.  At each such source point, $dv\ne0$.
The restriction of $\boldsymbol n$ from the smooth source branch to
the smooth image branch is locally an injective holomorphic map of
one complex variable.  Its first nonzero power must have exponent
one, since a higher power has more than one nearby preimage by
continuity of roots.  Its derivative is therefore nonzero.
The transverse Puiseux calculation now gives
$dh_0|_{\ker dv}\ne0$, so $d(v,h_0)$ is invertible.

Here $dv$ has no zero in a sufficiently small punctured source
neighbourhood.  To verify this from reducedness, a positive-dimensional
branch of the analytic set $\{dv=0\}$ would have $v$ constant, hence
zero, on its normalization.  It would therefore be contained in the
singular locus of the reduced plane curve $v=0$.  That singular locus
is isolated: otherwise a factor of the reduced equation would divide
both its partial derivatives, contrary to squarefreeness in
Weierstrass preparation.  This justifies choosing the initial ball
inside a representative with no other isolated critical points.

Fix $0<\rho<\rho_0$ and a small closed annulus around $|h_0|=\rho$.
The inverse image of that annulus at $v=0$ is compact and
$d(v,h_0)$ is invertible there.  Properness and a convergent-subsequence
argument therefore allow $\epsilon_0$ to be decreased so that this
invertibility holds on the annulus for all $|v|\le\epsilon_0$.
In particular all the vertical-boundary fibres are unramified, and
their cardinality is $S$, by the implicit function theorem and
properness, starting with their known cardinality at $v=0$.

At this point, before choosing the nonzero level, choose a smooth
central point $p$ with $0<|h_0(p)|<\rho/4$.  The implicit function
theorem for $v$ gives a section $\sigma(s)$ through $p$.  Fix a smaller
radius $\epsilon>0$ on which it is defined, with
$|h_0(\sigma(s))|<\rho/2$ on the closed disc of that radius.  Decrease
the preceding representative radius to this $\epsilon$ now.  This
order of choices ensures that a later choice of an allowed nonzero
level need not be changed in order to obtain a section.

\subsection{Capturing the singular points and choosing the final direction}

The critical set of $\boldsymbol n$ is a proper analytic subset of the
source: if all its $2\times2$ minors vanished identically, the local
rank theorem on the open rank-one locus would produce positive-dimensional
fibres, contradicting injectivity.  Its germ has only finitely many
curve branches.  The function $v$ is nonconstant on each such branch
$C$.  Indeed, on its normalization the map
$\boldsymbol n|_C$ is nonconstant and generically has nonzero derivative.
At those points the rank of $d\boldsymbol n$ is one and its image is
spanned by the tangent to $\boldsymbol n(C)$.  If $v|_C$ vanished
identically, $z_0$ would annihilate that image, giving $dv=0$ along
$C$, contrary to the isolated critical point of $v$.

On a normalized critical branch, therefore,
$v=t^q a(t)$ with $q\ge1$ and $a(0)\ne0$.  For every fixed small
parameter-disc annulus its absolute value has a positive minimum.
Consequently every intersection of these branches with $v=s$ tends to
zero as $s\to0$.  By the choice of the source ball above, these
are all the critical points of $\boldsymbol n$ in that representative.
In particular,
by choosing an allowed nonzero level sufficiently small, all of the
singular points listed in that level lie in $B_R$ and satisfy
$|h_0|<\rho/4$.  None is lost when the finite representative is used.
There are no additional singular points, since the restriction is
injective and is the normalization, as assumed in the question.

Fix such a level $s$.  If necessary replace
$h_0$ by $h=h_0+\lambda k$, with $\lambda$ arbitrarily small.  There
are only finitely many forbidden values of $\lambda$: those for which
the new projection vanishes on a central tangent or on the tangent of
one of the finitely many branches at the $Q_j$.  Every nonzero tangent
vector excludes at most one $\lambda$, since $(\xi_0,\zeta)$ are
coordinates.  Choose $\lambda$ outside that finite set.

All the representative properties can be retained with the already
chosen level $s$.  Explicitly, $k$ and $\det d(v,k)$ are bounded on
$\overline B_R$, while $\det d(v,h_0)$ has a positive lower bound on
the compact vertical-boundary annulus.  Choose $|\lambda|$ so small
that $h-h_0$ is smaller than its annular collar width and
$\lambda\det d(v,k)$ is smaller than half that lower bound.  Then
$d(v,h)$ is still invertible on $|h|=\rho$ for the same small range
of $v$.  Also require $|h-h_0|<\eta-\rho$ on
$\overline B_R$.  At a boundary point with $|v|\le\epsilon<\eta$
we have $|h_0|\ge\eta$, and hence $|h|>\rho$.
Thus no preimage of the closed target bidisc meets
$\partial B_R$.  Properness follows exactly as before.  Transversality
on the central branches gives a unique common zero of $(v,h)$:
their Puiseux units remain nonzero uniformly for sufficiently small
$\lambda$.  Compact fibres are finite, so the same openness and
connectedness arguments apply to the new $\Omega$.  The central sheet
count is still $S$.  This small change does not remove any $Q_j$ from
the representative.  Also require $|\lambda|$ small enough that the
already chosen section still satisfies $|h(\sigma(s))|<3\rho/4$ on
the closed base disc.  This is possible by compactness of its image.
We henceforth use $h$ and this representative, with the previously
fixed $\epsilon$ and the previously chosen nonzero level.

At every $Q_j$, transversality gives local degree exactly $a_j$.
We impose no condition that ordinary critical points be simple or that
their values be distinct.

\subsection{Connectedness of the compact level}

Distinguish the open set $\Omega$ from the set with its vertical boundary
adjoined:
\[
 E=\{p\in B_R:|v(p)|<\epsilon, |h(p)|\le\rho\},
 \qquad E_s=v^{-1}(s)\cap E.
\]
The fibres $E_s$ are compact.  For $s\ne0$ they are smooth complex
curves with smooth boundary, because $dv\ne0$ and $d(v,h)$ is
invertible near $|h|=\rho$.  The assertion of compactness here refers
to $E_s$, not to the open fibre in $\Omega$.

We give the needed local trivializations explicitly.  In source
coordinates $z_1,z_2$, put
\[
 Y_0=\frac{(\overline{v_{z_1}},\overline{v_{z_2}})}
 {|v_{z_1}|^2+|v_{z_2}|^2},\qquad
 Y_b=\frac{(h_{z_2},-h_{z_1})}
 {v_{z_1}h_{z_2}-v_{z_2}h_{z_1}}.
\]
The first is defined away from the zero level, and the second is
defined on the vertical-boundary collar.  Both satisfy $dv(Y)=1$,
and $dh(Y_b)=0$.  Choose a smooth function of $|h|$ equal to one near
$\rho$ and zero outside that collar, and use it to interpolate between
$Y_0$ and $Y_b$.  The resulting vector field $Y$ satisfies $dv(Y)=1$
and is tangent to the vertical boundary.  For a short smooth base path
$\gamma$, integrate
\[
 \dot p(t)=\gamma'(t)Y(p(t)).
\]
On a compact strip over a closed subdisc avoiding zero this vector field
and its derivatives are bounded.  Its solutions therefore exist for
the required sufficiently short interval; they cannot leave through the
vertical boundary, and reversing the path gives the inverse.  These
flows identify nearby compact fibres and their boundaries.  A finite
subdivision handles any compact base path in the punctured disc.
In particular the finite sets of their connected components are locally
constant and have unique transport along paths.

The open set $\Omega\setminus\{v=0\}$ is path connected.  Indeed
$\Omega$ is connected.  Connect two points by a chain of convex balls
inside it, choosing intermediate points outside $v=0$.  On a complex
line through two such points the restriction of $v$ is not identically
zero.  It has only finitely many zeros near the compact segment, and
small arcs in that complex line detour around them.  Performing this
operation in the finitely many balls gives a path in the complement.

The section $\sigma(s)$ was chosen before the nonzero level and before
the small projection perturbation.  By the explicit order of choices
above it is defined on the \emph{entire} disc $\Delta_\epsilon$ and
stays in $\Omega$ for the final projection.  No further shrinking,
and no change of the selected nonzero level, is required here.

Now fix a nonzero base value $s_*$.  Any component of $E_{s_*}$ contains
an interior point.  Join that point to $\sigma(s_*)$ by a path in
$\Omega\setminus\{v=0\}$.  Its projection is a loop in
$\Delta_\epsilon^*$.  Let $C_x$ be the component of the chosen initial
point, let $C_0$ contain $\sigma(s_*)$, and let $T$ be the bijection of
component sets supplied by transport around this loop.  The chosen
path proves $T(C_x)=C_0$.  The path obtained by applying the section
$\sigma$ to the same base loop proves $T(C_0)=C_0$.  Injectivity of
$T$ gives $C_x=C_0$; no equality of their initial components has been
assumed.  This proves that $E_{s_*}$ is
connected.  The argument uses a section on the entire punctured disc,
so no classification of its fundamental group or winding-number theorem
is needed.

On each compact smooth level, $h$ is nonconstant on each component
because its fibres are finite.  Its critical points are therefore
isolated and finite; there are none on the boundary collar.  Away from
their finitely many values, properness and the implicit function theorem
give a covering.  The punctured base disc is path connected, and its
regular fibres have the same cardinality as those near its outer circle,
namely $S$.

\subsection{Local factors, without a generic-Morse projection}

Fix a critical value $c$.  Its finitely many preimages have disjoint
normalization neighbourhoods and distinct $\zeta$ coordinates, by
injectivity of $\boldsymbol n$.  There are only finitely many such
fibres.  Rotate the complementary $\zeta$ coordinate by a fixed unit
complex number avoiding the finitely many directions for which two
distinct values in one of those fibres have the same real part.
Choose small $\zeta$ discs whose real projections are disjoint within
each such fibre.  Thus the sheets of each local cluster occupy
consecutive positions in their real-coordinate order.  Choose a sufficiently small
disc about $c$ so that each moving cluster stays in its own disc.  The
remaining sheets stay away from those discs.

At a singular point $Q_j$, the supplied branch-word construction in the
transverse coordinates $(h,\zeta)$ gives the positive event word
$\beta_j$ on $a_j$ sheets, with the cycle, generator-occurrence, and
length properties stated in the question.  We use that explicit
event-word construction, not only a statement about the knot type of
its closure.

At a smooth critical point, a normalization parameter can be chosen so
that $h-c=t^d$ for some $d\ge2$: take a holomorphic $d$th root of the
unit in its leading expansion.  Since the plane curve is smooth there,
$\zeta(t)$ has nonzero derivative and is injective on a smaller
parameter disc.  Replace $t^d$ locally by
\[
 p_u(t)=t^d-u\varepsilon t,\qquad 0\le u\le1,
\]
where $\varepsilon\ne0$ is sufficiently small.  On a fixed sufficiently
small circle $|h-c|=r$, all $d$ roots remain in that parameter disc and
remain distinct throughout the family.  To verify this explicitly,
the critical points satisfy $dt^{d-1}=u\varepsilon$, and their critical
values are $(1-d)t^d$, whose absolute values are less than $r$ when
$\varepsilon$ is small.  Thus no critical value reaches the boundary
circle; continuity of roots and injectivity of $\zeta(t)$ give a
continuous family of configurations of distinct points there.

For $u=1$, the $d-1$ critical points are simple and their critical
values are distinct: multiplication by $\varepsilon/d$ identifies the
distinct roots of $dt^{d-1}=\varepsilon$ with their $d$th powers.
Near each of them the projection is, after a holomorphic coordinate
change, $t\mapsto t^2$.  Its two roots make one counterclockwise half
twist.  Cutting its local base disc along paths to these $d-1$ values
therefore expresses its boundary configuration as a product of
$d-1$ conjugated positive single-generator events.  The preceding
boundary family identifies this factorization with the original
ordinary ramification.  The local conjugating words are simply
concatenated with its transport word in the global base.

If several clusters have the same original critical value, their motions
take place in disjoint $\zeta$ discs.  They can be performed consecutively
instead of simultaneously: replace their common angular parameter by
nondecreasing functions which run one cluster at a time.  Linear
interpolation of these parameter functions is an explicit homotopy of
configurations, because points in different discs never meet.  Thus the
same factorization separates all the singular-cluster factors and all
the ordinary positive single-generator factors, without assuming
distinct original critical values.

\subsection{Why the configuration homotopies give finite word relations}

We record the elementary finite-event verification used above and
below.  A continuous family of configurations of $S$ distinct complex
points on a compact parameter rectangle has a positive minimum pairwise
separation.  Its labels can be continued on the rectangle: cover the
image locally by disjoint root discs and continue labels across a fine
triangulation; consistency around each small triangle follows because
each root stays in its own disc.  Triangle deletion gives consistency
on the whole rectangle.

For a transport path which has not already been assigned an event word,
use a generic piecewise-affine approximation within these root discs,
keeping its endpoint configurations fixed.  The finitely many endpoint
configurations can be chosen with distinct real coordinates, by a small
rotation and by moving the cut endpoints on their regular circles.
The resulting path has finitely many transverse events and defines its
word.  Thus no finiteness or transversality assertion about an arbitrary
unperturbed continuous transport path is needed.

Approximate the labelled coordinates by piecewise-affine functions on a
sufficiently fine triangulation, keeping the error smaller than one
quarter of the minimum separation.  The approximation, and the straight
homotopy to it, still have distinct points.  On boundary paths already
used to define words, choose vertices between their finitely many
transverse events; sufficiently small changes preserve their words.
The vertex coordinates may now be chosen generically within those open
error bounds.  Only finitely many polynomial conditions are to be
avoided: a real-coordinate tie at a vertex, coincident zero lines in a
triangle, four equal real coordinates, three independent pair-equalities,
an event intersection on a mesh edge, or equality of the movie-parameter
coordinates of two unrelated exceptional events.  Each is a proper
polynomial condition in the free vertex coordinates.  For example,
three independent affine pair-differences have a common zero only if
their three coefficient rows have zero determinant; its determinant is
not identically zero.  Structural equalities involving three indices
are retained, since the three pair-differences then satisfy their
necessary linear relation.  Moving mesh vertices slightly, if needed,
also avoids vertical interior edges.  Finitely many proper polynomial
zero sets cannot fill any open choice of the vertex coordinates.

In every triangle a pair-equality
$\operatorname{Re}(z_i-z_j)=0$ is now a straight segment.  These segments
form a finite arrangement.  As a vertical movie-parameter slice moves,
its event word changes only in the following ways.
First, at a local turn of one pair-equality arc along a mesh edge,
two crossings of the same pair are born or disappear; their signs are
opposite because the real difference crosses zero in opposite
directions while the imaginary difference keeps its sign.  Second,
two independent pair-equality arcs cross, interchanging events on
disjoint adjacent pairs, which is far commutation.  Third, three real
coordinates coincide.  Their imaginary coordinates remain distinct,
and their six possible height orders give the braid relation, its
inverse, or one of the four mixed-sign forms
\[
 aba^{-1}=b^{-1}ab,\quad
 ab^{-1}a^{-1}=b^{-1}a^{-1}b,\quad
 a^{-1}ba=bab^{-1},\quad
 a^{-1}b^{-1}a=ba^{-1}b^{-1}.
\]
Each displayed relation follows from $aba=bab$ by multiplication and
cancellation.  The arrangement has finitely many segments, turns, and
intersections, so this procedure supplies a finite list of these word
operations.  It neither invokes Reidemeister's theorem nor identifies a
configuration-space fundamental group with an Artin group.

This argument first applies to a homotopy of loops with fixed basepoint
configuration, for which no event reaches the two endpoint edges of the
parameter rectangle.  In a loop homotopy $L_s$ with moving basepoint
$b(s)$, replace $L_s$ by the based loop
\[
 b|_{[0,s]}*L_s*\overline{b|_{[0,s]}}.
\]
These loops form a continuous homotopy with basepoint $b(0)$.
The preceding argument then compares the initial event word with the
conjugate of the final event word by the transport word of $b$.
Equivalently, keeping an unbased loop cut records a crossing through
that cut as a cyclic word rotation.  This treats the boundary case as
well as every interior event of the arrangement.

\subsection{Global factorization, exponent count, and transitivity}

Choose polygonal cuts from a regular boundary basepoint to disjoint small
discs about the distinct critical values.  The cuts are disjoint except
at their basepoint and have a specified cyclic order there.  The cut
disc with the small discs removed has a finite triangulation.  Removing
its triangles successively deforms its outer boundary to the ordered
product of path--small-circle--reverse-path loops.  The roots are
distinct throughout.  Apply the finite-event verification above and
the local factorizations just proved.  It gives an equality by explicit
finite event relations of the form
\[
 w_{\rm out}\ \sim\ \prod_{l=1}^{N+B}\gamma_l b_l\gamma_l^{-1},
\]
where the $N$ singular factors are the positive words $\beta_j$ and the
other $B$ factors are positive single generators.  Here
$B$ is the sum of the ordinary ramification orders $d-1$; it need not
be the number of original ordinary critical points.

Over $|h|=\rho$ the source sheets vary by the implicit function theorem
without collision as $v$ moves from zero to the chosen level, and
$\zeta\circ\boldsymbol n$ separates them by injectivity.  Thus their
outer configurations form a homotopy of closed $S$-braids.  If its
basepoint configuration moves, include its transport path; the resulting
word comparison is a conjugacy, again furnished by the finite-event
verification.  The central branch construction supplied in the question
then compares this outer word with $\mathbf d_S\mathbf w$.  All the
relations and conjugacies preserve exponent sums.  Therefore
\begin{align*}
 B
 &=2\delta(W)+S-b-\sum_j(2\delta_j+a_j-1)\\
 &=S-\sum_j a_j+N-1+\mu(V).
\end{align*}
This computation uses only the supplied branch lengths and delta
identities; no argument principle or Riemann--Hurwitz calculation is
needed here.

Finally let $\mathcal V$ be the finite set of critical values of the
original finite projection.  Remove the \emph{entire} finite set
$h^{-1}(\mathcal V)$ from the connected smooth normalization $E_s$,
not just its ramification points.  A connected real surface remains
path connected after finitely many interior points are removed: take
a finite chartwise path and detour around those points in small
coordinate discs.  A path in this complement between two points of
one regular fibre projects to a loop avoiding $\mathcal V$.

The monodromy of that loop is generated by the small critical-value
loops.  Explicitly, approximate it polygonally away from the punctures,
make it transverse to the chosen cuts, and contract each portion in the
cut region by deleting triangles.  Its finitely many crossings of the
cuts record a word in the path--circle--reverse-path loops.  This gives
the required expression without a fundamental-group presentation.
Consequently the permutations of the original critical-value loops
act transitively on the $S$ sheets.  Every such permutation is a product
of the permutations of the separated factors above, including the
$d-1$ factors replacing an ordinary higher ramification.  Hence the
permutations of all $N+B$ displayed factors also generate a transitive
group, as required.

\end{proof}


\section{A rectangle proof of the native O-saddle pairing}

\begin{lemma}[Native O calculus and incidence pairing; new combinatorial]
\label{lem:O-calculus}
\formalname{nativeOSaddleIncidence}
For every labelled O diagram of size at least two, the differentiated
rectangle maps $\Phi_i$ below are polynomial chain maps of degree
$+1$, annihilating the canonical class on homology.  An adjacent O
exchange has the explicit diagonal degree-$-1$ maps $F,F'$ below,
with shifts $(0,-1/2)$.  For either of the two exchanged labels $i$,
$K_i=\Phi_iF+F\Phi_i$ is a polynomial degree-zero homotopy equivalence,
preserves the canonical class, and intertwines all marking
derivatives on homology.  For any sequence of O exchanges, with
intervening degree-zero polynomial canonical-class maps intertwining
the colour derivatives, the product of as many colour derivatives
as O exchanges has the determinant action
\eqref{eq:determinant-pairing}.  In particular a rooted-tree incidence
matrix gives coefficient one.
\end{lemma}
\begin{proof}

All coefficients in this section have characteristic two.  Label the $n$
$O$ markings and put $R=\mathbb F_2[U_1,\ldots,U_n]$.  The full
multivariable grid differential is
\[
 D_Gx=\sum_y\sum_{r\in\operatorname{Rect}^{\circ}(x,y)}
       \left(\prod_{h=1}^n U_h^{n_h(r)}\right)y,
 \qquad n_h(r)=\#(r\cap O_h)\in\{0,1\}.
\]
Its square is zero: nonannular index-two domains have two decompositions
into empty rectangles; the remaining domains are the horizontal and vertical
annuli, and each $O_h$ occurs in precisely one annulus of each kind, giving
$U_h+U_h=0$.  This argument counts rectangles containing $X$ markings.
Write $d_G$ for the specialization $U_h=T$ for every $h$.

\subsection{Derivative maps and a rectangle homotopy}

Define
\[
 \Phi_h^G=\left.\frac{\partial D_G}{\partial U_h}\right|_{U_1=\cdots=U_n=T},
 \qquad
 \Phi_h^G(x)=\sum_{y,r} n_h(r)T^{n_O(r)-1}y.
\]
Every nonzero summand has $n_O(r)\ge1$, so this is a polynomial map,
including when the original complex has been localized in $T$.  It has
Maslov degree $+1$ and first filtration shift zero.  Since
\[
 A(y)-n_O(r)+1-A(x)=1-n_X(r),
\]
its second filtration shift is at most $1$.  If $O_h$ and an $X$ marking
occupy the same square, the second shift is at most zero.

Differentiating $D_G^2=0$ proves
\[
 d_G\Phi_h^G+\Phi_h^Gd_G=0.
\]
The differential is linear in each individual $U_h$, so its coefficient
of $U_h^2$ in $D_G^2$ proves $(\Phi_h^G)^2=0$ strictly.  For $h\ne k$,
twice differentiating proves
\[
 \Phi_h^G\Phi_k^G+\Phi_k^G\Phi_h^G
 =d_G B_{hk}+B_{hk}d_G,
 \qquad
 B_{hk}=\left.\frac{\partial^2D_G}{\partial U_h\partial U_k}
         \right|_{U=T}.
\]
Thus the induced derivative operators commute on homology.

Let $p$ lie in the interior of a vertical grid circle, away from the
horizontal grid circles, and suppose the two columns adjacent to this
circle contain $O_i$ and $O_j$.  For a rectangle $r$ let $a_p(r)$ be one
if its vertical boundary passes through $p$, and zero otherwise.  Define
the multivariable map
\[
 \mathcal A_p^G(x)=\sum_{y,r} a_p(r)
                      \prod_hU_h^{n_h(r)}y.
\]
Then
\begin{equation}\label{eq:relative-rectangle-action}
 D_G\mathcal A_p^G+\mathcal A_p^GD_G=(U_i+U_j)\operatorname{id}.
\end{equation}
Here is the entire additional check beyond the usual proof that $D_G^2=0$.
For a decomposition $r_1*r_2$ the contribution to the left side is weighted
by $a_p(r_1)+a_p(r_2)$.  This is the multiplicity at $p$, modulo two, of
the vertical boundary of the composite domain, and hence is independent
of its decomposition.  The two decompositions of every nonannular domain
therefore cancel with the same weight.  The vertical boundary of a
horizontal annulus is zero.  A vertical annulus has weight one precisely
when it is one of the two annuli adjacent to the circle containing $p$.
These two annuli contain $O_i$ and $O_j$, respectively.  Each has the usual
unique empty-rectangle decomposition at each state, giving exactly the
right side of \eqref{eq:relative-rectangle-action}.

Let $A_p^G=\mathcal A_p^G|_{U=T}$.  Differentiating
\eqref{eq:relative-rectangle-action} in $U_i$ and then specializing gives
\begin{equation}\label{eq:clifford-relation}
 \Phi_i^GA_p^G+A_p^G\Phi_i^G
 =\operatorname{id}+d_GB_i^G+B_i^Gd_G,
 \qquad
 B_i^G=\left.\frac{\partial\mathcal A_p^G}{\partial U_i}\right|_{U=T}.
\end{equation}
The same identity holds with $j$ in place of $i$.

\subsection{The O switch and its two detecting operators}

Consider adjacent nonwrapping columns $a,a+1$.  In $G$ suppose
\[
 O_i=(a+\tfrac12,r+\tfrac12),\qquad
 O_j=(a+\tfrac32,s+\tfrac12),\qquad r<s.
\]
In $H$ interchange these two $O$ markings, preserving their labels and
their rows.  All other $O$ markings and all $X$ markings stay fixed.
Set
\[
 e(x)=\mathbf1_{\{r<x(a+1)\le s\}},\qquad
 F(x)=T^{e(x)}x,\qquad F'(x)=T^{1-e(x)}x.
\]
The elementary $J$-function calculation gives
\[
 M_O^H(x)-M_O^G(x)=2e(x)-1.
\]
The rectangle Maslov formula consequently gives, for every empty rectangle
$q:x\longrightarrow y$,
\begin{equation}\label{eq:O-exponent-change}
 n_O^H(q)-n_O^G(q)=e(y)-e(x).
\end{equation}
It follows term by term that $F$ and $F'$ are chain maps of Maslov degree
$-1$, with shifts $(0,-\tfrac12)$, and
\[
 F'F=T\operatorname{id},\qquad FF'=T\operatorname{id}.
\]
The reverse local pattern is covered by interchanging $F,F'$ and $G,H$.
A cyclic permutation of columns reduces a wrapping column pair to this
case.

Define the degree-zero maps
\[
 K_i=\Phi_i^HF+F\Phi_i^G,\qquad
 K'_i=\Phi_i^GF'+F'\Phi_i^H.
\]
Both preserve the polynomial subcomplex, and both have shifts
$(0,\tfrac12)$.  Take $p=(a+1,r+\tfrac12)$.  A rectangle contains exactly
one of $O_i^G,O_i^H$ precisely when its vertical boundary passes through
$p$.  This is because the two points lie in adjacent squares on the same
horizontal level, and the unique grid circle between them is $x=a+1$.
Using \eqref{eq:O-exponent-change} to compare every individual rectangle
therefore proves the exact identities
\begin{equation}\label{eq:conjugated-derivative}
 F'\Phi_i^HF=T\Phi_i^G+A_p^G,
 \qquad F\Phi_i^GF'=T\Phi_i^H+A_p^H.
\end{equation}
For example, the coefficient of a rectangle $q:x\to y$ on the left of
the first identity is
\[
 n_i^H(q)T^{e(x)+n_O^H(q)-1+1-e(y)}
 =n_i^H(q)T^{n_O^G(q)}.
\]
The right side has coefficient
$(n_i^G(q)+a_p(q))T^{n_O^G(q)}$, which is the same.

Expanding and using $(\Phi_i^G)^2=(\Phi_i^H)^2=0$ gives
\[
 \begin{split}
 K'_iK_i
 &=\Phi_i^G(F'\Phi_i^HF)+(F'\Phi_i^HF)\Phi_i^G\\
 &=\Phi_i^GA_p^G+A_p^G\Phi_i^G
 =\operatorname{id}+d_GB_i^G+B_i^Gd_G.
 \end{split}
\]
The identical calculation in $H$ gives
$K_iK'_i=\operatorname{id}+d_HB_i^H+B_i^Hd_H$.
Thus $K_i$ is a polynomial-preserving, degree-zero chain homotopy
equivalence, with the displayed explicit inverse and homotopies.

For later use, the induced $K_i$ intertwines every labelled derivative
operator.  To prove this, first work over the ring in which only
$U_i=U_j=u$ has been imposed.  The same rectangle calculation makes
$F=u^e\operatorname{id}$ a chain map over this ring.  Differentiation in
$u$ proves
\[
 K_i+K_j=d_HJ+Jd_G,\qquad J(x)=e(x)x.
\]
Differentiation in any other $U_k$ proves the strict identity
$\Phi_k^HF=F\Phi_k^G$.  Now $\Phi_i^HK_i=K_i\Phi_i^G$ follows strictly
from the square-zero identity.  For $j$ it follows on homology by
replacing $K_i$ with $K_j$.  For any other $k$, it follows from the strict
identity for $F$ and the commutation, on homology, of the derivative maps.
In particular,
\begin{equation}\label{eq:K-derivative-naturality}
 (\Phi_k^H)_*(K_i)_*=(K_i)_*(\Phi_k^G)_*
 \quad\hbox{for every label }k.
\end{equation}

\subsection{Why every derivative annihilates the canonical class}

We first prove the slightly stronger fact
\begin{equation}\label{eq:positive-polynomial-image-zero}
 \operatorname{im}\bigl(H_m(P(G))\longrightarrow H_m(C(G))\bigr)=0
 \qquad(m>0).
\end{equation}
This uses the explicitly proved $K_i$ equivalences and a finite
permutation inequality; it uses no homology computation or grid isotopy
theorem.  The differential, Maslov grading, and polynomial subcomplex
ignore $X$.  Work temporarily with $O$-only grids and bubble-sort the
$O$ permutation to $O(c)=n-1-c$ by adjacent column exchanges.  Each
exchange has the degree-zero polynomial-preserving equivalence $K_i$
constructed above.  Its proof did not use
\eqref{eq:positive-polynomial-image-zero} or any assertion about the
canonical class.  Each $K_i$ is injective on Laurent homology and maps
the polynomial image into the polynomial image.  It therefore suffices
to verify \eqref{eq:positive-polynomial-image-zero} for antidiagonal $O$.
If ordinary grid notation is desired throughout, choose an arbitrary
disjoint $X$ permutation at each stage; neither the complexes in this
argument nor their first filtrations change.

For a state $x$, put $\pi(c)=n-1-x(c)$, let
$F(\pi)=\sum_c|\pi(c)-c|$, let $D(\pi)=\#\{c:\pi(c)<c\}$, and write
$\operatorname{inv}(\pi)$ for its inversion count.  Directly from $J$,
\[
 M_O(x)=\operatorname{inv}(\pi)-\sum_c|\pi(c)-c+1|+1
        =\operatorname{inv}(\pi)-F(\pi)-n+2D(\pi)+1.
\]
If $c(\pi)$ denotes the number of cycles, then
\[
 \operatorname{inv}(\pi)+n-c(\pi)\le F(\pi),
 \qquad D(\pi)\le n-c(\pi).
\]
For completeness, prove the first inequality by repeatedly taking the
least nonfixed index $i$, putting $j=\pi^{-1}(i)>i$, and interchanging
the values at $i,j$.  Write $v=\pi(i)>i$.  This decreases the footrule
by $2\min(v-i,j-i)$ and decreases the inversion count by
\[
 1+2\#\{k:i<k<j,\ i<\pi(k)<v\}
 \le 2\min(v-i,j-i)-1.
\]
It therefore decreases $F-\operatorname{inv}$ by at least one.  The
transposition splits off the fixed point $i$ from its cycle, and hence
exactly $n-c(\pi)$ such steps produce the identity.  This proves the
first inequality.  The second follows because each cycle of length
$q$ has at most $q-1$ deficient indices.  Consequently
\[
 M_O(x)\le -\bigl(n-c(\pi)\bigr)-n+2D(\pi)+1
          \le D(\pi)-n+1\le0.
\]
Every polynomial monomial thus has nonpositive Maslov degree.  This
proves \eqref{eq:positive-polynomial-image-zero}.  Equality $M_O=0$ forces
$D(\pi)=n-1$ and hence
$\pi=(n-1,0,1,\ldots,n-2)$, giving the unique state
$\theta=(0,n-1,n-2,\ldots,1)$.  Thus the degree-zero polynomial
chain group is the one-dimensional space generated by $\theta$.
The supplied canonical-class theorem says that its degree-zero
homology has nonzero image in Laurent homology.  It follows at once
that $\theta$ is a cycle with nonzero Laurent class, and this class is
the canonical class in the normal form.

Since $\Phi_i^G$ is polynomial preserving of degree $+1$, applying it to
a polynomial degree-zero representative of the canonical class and using
\eqref{eq:positive-polynomial-image-zero} gives
\[
 (\Phi_i^G)_*\operatorname{can}(G)=0.
\]
Because $K_i$ has degree zero, preserves the polynomial subcomplex, and
is an isomorphism on homology, Lemma N gives
$(K_i)_*\operatorname{can}(G)=\operatorname{can}(H)$.  We have therefore
proved the explicit O-switch pairing identities
\begin{equation}\label{eq:O-switch-top-pairing}
 (\Phi_i^H)_*F_*\operatorname{can}(G)
 =(\Phi_j^H)_*F_*\operatorname{can}(G)=\operatorname{can}(H),
 \qquad
 (\Phi_k^H)_*F_*\operatorname{can}(G)=0\quad(k\notin\{i,j\}).
\end{equation}

\subsection{Component colours and the determinant formula}

The fresh-label hypothesis can be removed.  Assign each $O$ marking a
colour in a fixed finite set.  At a given stage put
\[
 \Psi_c=\sum_{\operatorname{col}(i)=c}\Phi_i.
\]
These are degree-one chain maps, polynomial preserving, whose induced
operators commute and square to zero.  Indeed, the squares of the
individual $\Phi_i$ vanish, and every mixed summand in $\Psi_c^2$ is an
anticommutator, already proved null-homotopic.  Each $\Psi_c$ annihilates
the canonical class by the positive-degree polynomial vanishing theorem.

Suppose an O switch $F_e$ exchanges labels $i_e,j_e$.  The colour is
attached to the labelled marking and moves with it.  Set
\[
 b_{c,e}=\mathbf1_{\{\operatorname{col}(i_e)=c\}}
          +\mathbf1_{\{\operatorname{col}(j_e)=c\}}
          \ \in\mathbb F_2.
\]
On homology the identities already proved give
\begin{equation}\label{eq:colour-switch-commutator}
 (\Psi_c^{\rm after})_* (F_e)_*
 =(F_e)_*(\Psi_c^{\rm before})_*+b_{c,e}(K_e)_*.
\end{equation}
The same $K_e$ can be used for either endpoint of the edge since
$K_{i_e}\simeq K_{j_e}$.  Moreover $K_e$ intertwines every $\Psi_c$.

Let $F_1,\ldots,F_d$ be any sequence of $d$ O switches, and choose
$d$ colours $c_1,\ldots,c_d$.  Let $B=(b_{c_r,e})_{1\le r,e\le d}$.
Then
\begin{equation}\label{eq:determinant-pairing}
 (\Psi_{c_1}^{\rm final})_*\cdots(\Psi_{c_d}^{\rm final})_*
 (F_d)_*\cdots(F_1)_*\operatorname{can}(G_0)
 =\det(B)\,\operatorname{can}(G_d).
\end{equation}
To prove the formula, move the derivatives successively to the right
using \eqref{eq:colour-switch-commutator}.  At a position still occupied
by $F_e$ a derivative either passes through or consumes that position,
replacing $F_e$ by $b_{c,e}K_e$.  A later derivative passes through a
$K_e$ without any extra term, so a position can be consumed at most once.
Every term in which a derivative reaches the initial canonical class
vanishes.  Since there are exactly $d$ derivatives and $d$ positions,
each surviving term consumes every position exactly once.  Such terms
are indexed by bijections between the derivatives and the positions.
Their coefficients sum to the permanent of $B$, which equals its
determinant in characteristic two.  In every such term the operator
product is exactly $K_d\cdots K_1$, in its original order; no
commutativity of a $K$ with an $F$ is being used.  That product sends the
initial canonical class to the final one.  This proves the formula.

There may also be arbitrary intervening degree-zero, polynomial-preserving
canonical-class maps $P$ which intertwine the colour derivatives.  The
same expansion passes derivatives through each $P$ and proves the
identical formula with the $P$'s left in their original positions.

Here is the precise coefficient rule that gives colour naturality for
native moves.  Replace the independent marking variables by colour
variables $V_c$, namely $U_i=V_{\operatorname{col}(i)}$, and denote the
resulting differential by $D^{\rm col}$.  For any chain map
$P^{\rm col}$ defined over this polynomial ring, differentiation of
\[
 D_H^{\rm col}P^{\rm col}=P^{\rm col}D_G^{\rm col}
\]
in $V_c$, followed by $V_c=T$, proves
\[
 \Psi_c^HP+P\Psi_c^G=d_H P_c+P_cd_G,
 \qquad P_c=\left.\frac{\partial P^{\rm col}}{\partial V_c}\right|_{V=T}.
\]
For a stabilization the new $O$ marking must inherit the colour of the
old $O$ marking whose variable it is identified with; a destabilization
uses the reverse rule.  Under that rule the native polygon formulas
lift to the colour ring: paired domains have the same complete marking
monomial, and the exceptional stabilization annulus terms have the two
equal colour variables.  This explains exactly why differentiating by
one arbitrarily chosen representative marking is insufficient, whereas
the full colour derivative is natural under these moves.

Finally, suppose the graph whose vertices are the colours and whose
edges are the O switches is a tree on $d+1$ vertices.  Choose one root
colour and choose all the other $d$ colours in
\eqref{eq:determinant-pairing}.  Then $B$ is its incidence matrix with
the root row removed, so $\det(B)=1$.  A proof requiring no graph
linear algebra is induction on $d$: a nonroot leaf exists, its row has
exactly one entry equal to one, and expansion along that row removes
that leaf and its incident edge.  The empty tree has determinant one.
The same argument applies component by component to a forest with one
root per component.

To apply this to deaths, colour every marking of each final isolated
unknot by the colour of that component.  Ordinary moves reducing that
component to a doubly marked square preserve its colour derivative by
the preceding rule.  At the doubly marked square the derivative has
second filtration shift zero.  The degree-one death can be written
$R\Psi_c$, where $R$ is the degree-zero top projection and has second
shift $+\tfrac12$.  The derivative for a different dying component
commutes past $R$ on homology, since that projection changes no marking
of the other colour.  All deaths therefore give a product of the
component derivatives followed by the degree-zero top projections.
If their switch incidence graph is the rooted forest just described,
\eqref{eq:determinant-pairing} proves that the complete map carries the
canonical class to the canonical class.

\end{proof}


\section{Auxiliary contractions and four-mark normalization}
\label{sec:aux-death}

In this section only, an \emph{auxiliary marking pair} is a pair of
permutations $(O,X)$ for which some values $O(c)=X(c)$ are allowed.  This
does not enlarge the class of native endpoints.  Its chain module,
differential, and two grading functions are given by exactly the formulas
in the question.  The complex and its polynomial image depend only on
$O$.  Consequently its canonical class is defined by identifying this
complex with the native complex of any valid pair having the same $O$.
Such a valid pair exists in every size at least two; for example, take
$X(c)=O(c+1)$ cyclically.  The resulting class is independent of this
choice, because these are literal identifications of the chain complex
and of its polynomial subcomplex.


The derivative $\Phi_i$ is the explicitly defined rectangle map of
Lemma~\ref{lem:O-calculus}.  At a coincident O/X square it has
degree $+1$ and shifts $(0,0)$, by
$A(y)-n_O(q)+1-A(x)=1-n_X(q)\le0$.

\subsection{Exact local coordinates}

\begin{lemma}[Local auxiliary contraction; new combinatorial]\label{lem:aux-death}
\formalname{nativeAuxiliaryContraction}
Let $G$ be a valid grid of size $N\ge2$ and choose an X square
$(c,r)$, so $G.X(c)=r$.  Insert a column immediately to the right of
column $c$ and a row immediately above row $r$.  Embed the old indices by
\[
 \iota_c(a)=\begin{cases}a&a\le c,\\a+1&a>c,\end{cases}
 \qquad
 \iota_r(b)=\begin{cases}b&b\le r,\\b+1&b>r.\end{cases}
\]
Let $B$ be the auxiliary pair retaining all embedded old markings and
having an additional O and X together in square $(c+1,r+1)$.  Label this
new O marking $i$.  There is an explicitly specified chain map
\[
 D_B:C(B)\longrightarrow C(G)
\]
of degree $+1$ and shifts $(0,\tfrac12)$.  There is also a degree-zero
map $\rho_B:C(B)\to C(G)$ with those shifts and with
$\rho_{B*}\operatorname{can}(B)=\operatorname{can}(G)$, such that
\[
 D_B=\rho_B\Phi_i.
\]
\end{lemma}

\begin{proof}
Let $A$ have the same O markings as $B$, and replace the two X markings
$(c,r),(c+1,r+1)$ by
$(c,r+1),(c+1,r)$.  In the indicated $2\times2$ block $A$ therefore has
\[
 \begin{array}{c|cc}
 &c&c+1\\ \hline
 r+1&X&O_i\\
 r&\varnothing&X
 \end{array}
\]
and all its other markings are the embedded markings of $G$.  This is
exactly an X stabilization of $G$, with the southwest square empty.
It is valid: the old O in column $c$ was not in row $r$, so after
insertion it occupies neither of the displayed X squares.  All remaining
validity conditions follow from those of $G$.

Use the explicit type-R destabilization count
$R_A:C(A)\to C(G)$ defined below.  Section~\ref{sec:type-r-independent}
proves directly that this particular formula has degree zero, shifts
$(0,0)$, and preserves the canonical class.  The identity of the
underlying O-complexes defines
$I:C(B)\to C(A)$.  The usual adjacent-X grading calculation depends
only on the X permutations, not on disjointness from O, and therefore
applies here: $M_X^A(x)-M_X^B(x)\in\{-1,1\}$ for every state.  Hence
$I$ has shifts $(0,1/2)$.  Because it is a literal O-complex and
polynomial-subcomplex identity, it preserves the canonical class.

Set
\[
 \rho_B=R_A I,\qquad D_B=R_A I\Phi_i.
\]
The preceding derivative calculation applies to $B$, since $O_i$ and
its X coincide.  Thus $\Phi_i$ has degree $+1$ and shifts $(0,0)$,
which proves every stated chain, degree, and filtration property.
The canonical-class identity for $\rho_B$ follows from the identities
for $I$ and $R_A$.
\end{proof}

\begin{corollary}[Paired contraction; new combinatorial]\label{cor:paired-aux-contraction}
\formalname{pairedAuxiliaryContraction}
Suppose an O-switch map $F:C(K)\to C(B)$ has degree $-1$, shifts
$(0,-1/2)$, and satisfies
\[
 \Phi_{i*}F_*\operatorname{can}(K)=\operatorname{can}(B).
\]
Then $D_BF$ has degree zero, shifts $(0,0)$, and sends
$\operatorname{can}(K)$ to $\operatorname{can}(G)$.
\end{corollary}
\begin{proof}
Compose the maps and their grading bounds, and use
$D_B=\rho_B\Phi_i$ and the canonical-class identity for $\rho_B$.
\end{proof}

\subsection{Labelled naturality and the persistent root variable}

For global compositions, choose a label $0$ which survives all the
contractions.  Place the auxiliary square in Lemma~\ref{lem:aux-death}
next to the X in the row of $O_0$.  The O marking traditionally paired
with $O_i$ in this stabilization is then exactly $O_0$.
Use the coefficient ring
\[
 \mathbb F_2[T,U_j:j\ne0,i],\qquad U_0=U_i=T,
\]
for its top destabilization map.  The explicit polynomial representative
of that map is the following.  Put $p=(c+1,r+1)$ at the crossing of
the newly inserted circles, and let $e(y)$ delete $p$ from a state
$y$ containing it, followed by removal of the new row and column.
A domain $q:x\to y$ is an array of integers $q_{a,b}$ on the
grid squares, with cyclic indices, satisfying at every intersection
$(a,b)$ the boundary equation
\[
 q_{a,b}+q_{a-1,b-1}-q_{a-1,b}-q_{a,b-1}
   =\mathbf 1_{\{x(a)=b\}}-\mathbf 1_{\{y(a)=b\}}.
\]
This is the northeast-plus-southwest minus northwest-minus-southeast
corner convention used for the empty rectangles.  Count nonnegative
domains from $x$ to $y$ of type R: at every point of
$(x\cup y)\setminus\{p\}$ at least three adjacent multiplicities
vanish; at $p$ three
multiplicities are $k$ and the southeast multiplicity is $k+1$; and
$y\setminus x$ has $2k+2$ points.  A domain contributes
\[
 \left(\prod_{j\ne i}U_j^{n_{O_j}(q)}\right)e(y).
\]
The set of domains is finite: the number $2k+2$ is at most the grid size,
and the zero-multiplicity conditions determine a domain of each boundary
type uniquely.  There is no factor $U_i$ in this weight.  This is the
usual top projection representative of the native X-destabilization map.

The verification of its chain equation retains these separate
variables.  Its cancellations have two forms.  In the first, a rectangle
and a type-R domain are cut in the other order, with the same total
domain; the marking monomials are therefore identical.  In the second,
the alternative domains differ by one of the thin annuli in the inserted
row or column.  These annuli contain $O_i$ or the O marking in the row
of the lower X, namely $O_0$, and no other O markings.  Their coefficient
contributions agree after $U_i=U_0=T$.  No identification of any other
$U_j$ is involved.  Equivalently, before that identification the only
uncancelled term is $(U_i+U_0)$ times the complementary type-L domain
count.  This establishes the required chain equation over the displayed
ring, rather than merely after setting every variable to $T$.

Consequently, for each surviving label $j\ne0,i$, differentiating the
chain equation and then setting all variables to $T$ yields
\[
 \Phi_j^G R_A+R_A\Phi_j^A
 =\partial R_{A,j}+R_{A,j}\partial,
 \qquad
 R_{A,j}=\left.\frac{\partial R_A}{\partial U_j}\right|_{U=T}.
\]
The identity $I$ changes no O label or differential, so the same
homotopy relation holds for $\rho_B$.  This is exactly the naturality
needed to move the derivative operators of the remaining contractions
past $\rho_B$.  The root label is never differentiated.  If instead a
surviving label is used as the variable paired with $U_i$, differentiation
also produces $R_A\Phi_i$; choosing the persistent root as above avoids
that extra term.

\subsection{Moving and removing an arbitrary separated auxiliary component}

The preceding construction also works when $G$ has other doubly-marked
squares, provided the chosen X square $(c,r)$ is not itself doubly
marked.  The type-R formula depends on the O markings and the
position of the stabilization, but not on any other X marking.
Its chain equation and its independent canonical-class proof in
Section~\ref{sec:type-r-independent} hold for all O diagrams.
The Alexander estimate there is the negative of a sum of
nonnegative X multiplicities.  It therefore remains valid when some
of those X markings coincide with O markings.  The derivative estimate
at the distinguished doubled square and the adjacent-X-exchange
estimate likewise remain valid.  Hence Lemma~\ref{lem:aux-death}
applies in this more general situation.

For the same reason, commutation maps extend to auxiliary pairs whenever
their marking segments are noninterleaving.  More explicitly, use the
usual empty-pentagon count, putting the O-multiplicity monomial on every
pentagon.  Its differential and homotopy-inverse verifications pair
domains with identical O multiplicities, exactly as for a native
commutation.  None of these verifications uses disjointness of the X and
O markings.  The grading comparison for a pentagon is unchanged and its
Alexander loss is the total nonnegative X multiplicity.  Thus both
directions have degree zero and shifts $(0,0)$, and the homotopy
equivalence preserves the polynomial image and hence the canonical
class.  Keeping separate O variables in these domain identities also
gives labelled derivative naturality by differentiation.

Here is an explicit check of the degenerate segment used below.  Suppose
the doubled square is in the left of the two adjacent columns, whose
horizontal coordinates are $(a,a+1)$ and $(a+1,a+2)$, and its marking
height is $v+1/2$.  Choose a small open interval $J$ about $v+1/2$,
contained in that row and containing no marking height of the other
column.  Replace the intervening vertical circle $x=a+1$ by the graph
$x=f(y)$, periodic in $y$, with $f=a+1/4$ on a smaller interval about
$v+1/2$, $f=a+7/4$ outside $J$, and exactly one monotone transition at
each end of $J$.  This graph avoids every marking, meets each horizontal
grid circle once, and intersects $x=a+1$ exactly twice.  It puts the
doubled marking in the right region and both markings of the other
column in the left region, and hence gives precisely the commuted
diagram after straightening.  Reflection gives the construction for a
doubled right column.  Interchange the coordinates for a row move.

To check every grading and cancellation at the coinciding punctures,
choose a disk around the doubled marking disjoint from all grid circles
and from this new graph.  Temporarily separate O and X by arbitrarily
small distinct displacements inside that disk.  The small segment between
the displaced heights still lies in $J$, so the same commutation curves
work.  Each rectangle, pentagon, and hexagon domain has constant
multiplicity on the disk; therefore its O and X multiplicities are
unchanged by bringing the two points back together.  The pairwise orders
of the X markings and their orders relative to every state coordinate
are unchanged as well, so the $J$-grading comparisons are unchanged.
The rectangle--pentagon and pentagon--hexagon cancellation proofs use
only these incidences and multiplicities.  They thus verify the maps,
homotopies, and filtration bounds also at the coinciding marking.  This
is a direct verification of the polygon formulas on auxiliary pairs,
not an appeal to a topological invariance theorem for them.

\begin{lemma}[Singleton transport; new combinatorial]\label{lem:dummy-transport}
\formalname{auxiliarySingletonTransport}
Let $E$ be an auxiliary pair of size $n\ge3$ with
$E.O(u)=E.X(u)=v$, and label this doubled marking $i$.  Delete its column
and row and call the resulting auxiliary pair $G$, of size $n-1$.
Choose a surviving root marking whose row $r$ in $G$ is not doubly
marked, and put $c=G.X^{-1}(r)$.  There is an explicit finite sequence of
noninterleaving commutations taking $E$ to $B(G,c,r)$, with all labels
retained.  In particular there is a degree-$+1$ map
\[
 \mathcal E_i:C(E)\longrightarrow C(G)
\]
with shifts $(0,1/2)$, and a degree-zero canonical-class-preserving map
$\rho_i:C(E)\to C(G)$ with the same shifts, such that
\[
 (\mathcal E_i)_*=(\rho_i)_*(\Phi_i)_*.
\]
The maps intertwine every surviving derivative except possibly that of
the root; the root derivative is not needed.
\end{lemma}

\begin{proof}
For clarity, deleting the dummy means the exact formulas
\[
 G.Y\bigl(a-\mathbf1_{a>u}\bigr)
 =E.Y(a)-\mathbf1_{E.Y(a)>v}
 \quad(a\ne u,\quad Y\in\{O,X\}).
\]
The desired insertion indices are $u_*=c+1$ and $v_*=r+1$.
If $u<u_*$, commute the dummy column successively across the pairs
$(u,u+1),(u+1,u+2),\ldots,(u_*-1,u_*)$; if $u>u_*$, use the reverse
list $(u-1,u),(u-2,u-1),\ldots,(u_*,u_*+1)$.  Use no column moves if
$u=u_*$.  Then apply the analogous list of row commutations to move
the dummy row from $v$ to $v_*$.  There are exactly
$|u-u_*|+|v-v_*|$ commutations.

Every one is noninterleaving.  The dummy column has both endpoints at
the same row, and no other column has either marking in that row,
because the dummy already uses both the unique O and the unique X of
that row.  A singleton segment cannot interleave with a segment whose
endpoints are distinct from that singleton.  The same argument applies
to the dummy row.  It applies also when the adjacent column or row
belongs to another doubled marking.  These moves never change the
relative order of the retained rows or the retained columns.  Therefore,
after deletion of the dummy, their result is still exactly $G$ with its
original indexing.  With insertion positions $c+1,r+1$, the resulting
pair is exactly $B(G,c,r)$, as asserted.

Let $P:C(E)\to C(B(G,c,r))$ be the composition of the commutation maps.
Set
\[
 \rho_i=\rho_B P,
 \qquad
 \mathcal E_i=D_B P=\rho_B\Phi_i^B P.
\]
All the degree, filtration, and canonical-class assertions follow from
the preceding lemmas.  Differentiating the fully labelled commutation
identities proves $\Phi_i^B P\simeq P\Phi_i^E$, giving the displayed
identity on homology.  Naturality for the other surviving labels follows
from the same commutation identities and from the root-variable
destabilization identity proved above.
\end{proof}

Consequently, after any prior algebraic or geometric cancellation has
left a list of doubled singleton components, their removal is an
explicit operation on the marking permutations: apply the adjacent
transposition lists in Lemma~\ref{lem:dummy-transport} and then the
displayed contraction.  There is no additional link-equivalence or
component-placement assertion at this stage.  The remaining hypothesis
for a global application is that its previous construction has indeed
produced these singleton components and a surviving root row with
distinct O and X markings.

\subsection{The component-colour version}

The same statements hold when several O markings have one common
component colour.  Assign a variable $V_\alpha$ to each colour and
substitute $U_j=V_{\operatorname{colour}(j)}$ in the multivariable
differential.  Its derivative, after collapse to $T$, is
\[
 \Phi_\alpha=\sum_{j:\operatorname{colour}(j)=\alpha}\Phi_j.
\]
This follows by the polynomial chain rule, including terms whose
rectangle contains several markings of that colour.  Thus derivative
naturality follows directly from the same chain-map identities, with no
requirement that a colour have only one O marking.  In an ordinary
stabilization, give its new O the colour of the component on which the
stabilization is performed.  The two variables identified by its top
projection then already coincide.

In a dummy contraction, let $\delta$ be the colour of the deleted
singleton and $\alpha_0$ the persistent root colour.  The top projection
substitutes $V_\delta=V_{\alpha_0}$.  Differentiation in a retained colour
different from $\alpha_0$ gives ordinary derivative naturality;
differentiation in $V_{\alpha_0}$ produces
\[
 (\Phi_{\alpha_0}^{\rm out})_*\rho_*
 =\rho_*\bigl((\Phi_{\alpha_0}^{\rm in})_*
                    +(\Phi_\delta^{\rm in})_*\bigr).
\]
The contraction itself is $\rho_*\Phi_{\delta*}$.  Its extra term
vanishes because $\Phi_{\delta*}^2=0$.  More generally every colour
derivative squares to zero on homology: expand its square as a sum of
individual squares and mixed anticommutators, use $\Phi_j^2=0$, and
differentiate $\mathcal D^2=0$ in two different variables to see that
each mixed anticommutator is null-homotopic.  This verifies the
substitution rule required when contracting several dummy components
in succession.

\subsection{Collapsing the four-mark unknot block}

\begin{proposition}[Four-mark collapse; new combinatorial]\label{prop:four-mark-collapse}
\formalname{fourMarkBlockCollapse}
Suppose an auxiliary grid has a doubled singleton in square $(c,r)$,
and has, elsewhere, a persistent component with a row whose O and X
are distinct.  Expand the singleton by inserting a column immediately
to its right and a row immediately above it, and replace it by the
four-mark block
\[
 \begin{array}{c|cc}
 &c&c+1\\ \hline
 r+1&X&O_i\\
 r&O_j&X.
 \end{array}
\]
Here $j$ is the old O label and $i$ is the newly introduced label.
There is an explicit map from this expanded grid to the original
auxiliary grid of degree zero and shifts $(0,0)$, preserving the
canonical class.  It intertwines the component-colour derivatives when
$i$ and $j$ are assigned the original singleton's colour.
\end{proposition}

\begin{proof}
This four-mark pattern is not an ordinary three-corner
destabilization, and we do not treat it as one.  Instead, exchange its
two O markings.  The explicit O-switch map $F$ has degree $-1$ and
shifts $(0,-1/2)$.  Its target has two doubled singleton squares, at
$(c,r+1)$ and $(c+1,r)$.  The former carries the new label $i$, and
the latter carries the old label $j$, because labels travel with their
original O rows under this exchange.

Apply Lemma~\ref{lem:dummy-transport} to contract the singleton $i$,
using the persistent component as root.  Call that degree-$+1$ map
$\mathcal E_i$.  The composite $\mathcal E_i F$ has degree zero and
shifts $(0,0)$.  The O-switch pairing identity gives
\[
 (\Phi_i)_*F_*\operatorname{can}
 =\operatorname{can}
\]
in the intermediate auxiliary complex.  Since
$(\mathcal E_i)_*=(\rho_i)_*(\Phi_i)_*$ and $\rho_i$ preserves the
canonical class, the composite preserves it as well.

Deleting the singleton $i$ means deleting column $c$ and row $r+1$.
The retained singleton at $(c+1,r)$ therefore becomes $(c,r)$.
Every other inserted-row or inserted-column shift is undone by these
deletions.  The exact target marking permutations are consequently
those of the original auxiliary grid.

Here is also the colour-naturality check; it is needed because the
source has two O markings of the singleton's colour.  Write $A$ for
the expanded source and $B$ for the intermediate two-singleton pair.
On homology put
\[
 K_i=\Phi_i^B F+F\Phi_i^A,
 \qquad K_j=\Phi_j^B F+F\Phi_j^A.
\]
The O-switch calculation gives $K_i=K_j$ on homology, and $K_i$
intertwines each derivative.  Using commutation of derivatives on
homology and their square-zero identities, we obtain
\begin{align*}
 \Phi_j^B\Phi_i^B F
       +\Phi_i^B F(\Phi_i^A+\Phi_j^A)
 &=\Phi_i^B K_j+\Phi_i^B F\Phi_i^A\\
 &=\Phi_i^B K_i+K_i\Phi_i^A=0.
\end{align*}
After applying $\rho_i$, this is precisely naturality between the
surviving $\Phi_j$ and the source colour derivative
$\Phi_i+\Phi_j$.  Other markings of that colour commute through the
construction.  Derivatives of all other retained colours are natural
by the preceding substitution calculation; for the persistent root,
its possible extra term vanishes by $\Phi_i^2=0$.
\end{proof}


\section{The explicit polynomial stabilization maps}
\label{sec:type-r-independent}

\begin{lemma}[Polynomial top domain maps; new combinatorial]
\label{lem:top-domain-maps}
\formalname{nativeTopDomainMaps}
For the X:SW stabilization and the incoming type-R formula specified
in Section~\ref{sec:aux-death}, the map $R$ has degree zero, shifts
$(0,0)$, and canonical-class identity.  Its multivariable chain
defect is precisely $(U_i+U_0)L$, where $i$ is the new O and $0$
the old O in the stabilized X row.  The outgoing formula $S$ below
is a polynomial chain map of degree zero and shifts $(0,0)$,
preserving the canonical class.  Both maps are natural on homology
for component-colour derivatives after assigning the new O its
component colour.  Together with the supplied commutation maps and
corner reduction this proves that naturality for every ordinary
move in both directions.
\end{lemma}
\begin{proof}

The map used in the auxiliary contraction is the explicit type-R count,
not an arbitrarily chosen projection supplied by the grid-move theorem.
We verify its grading and its canonical-class action directly.  Write
$G$ for the old grid, $A$ for its displayed X:SW stabilization, $i$ for
the new O, and $0$ for the old O in the row of the stabilized X.  The
definition and the domain cancellation give, before specialization,
\begin{equation}\label{eq:R-full-variable-identity}
 D_G R+R D_A=(U_i+U_0)L.
\end{equation}
Here $L$ counts the trivial domain, together with the nontrivial
positive domains with southwest multiplicity $k-1$, the other three
multiplicities $k$, and $2k+1$ moving state coordinates.  Its monomial,
like that of $R$, omits $U_i$.  In particular $R$ is a chain map when
only $U_i=U_0$ is imposed; all other variables remain independent.

We use the following elementary facts about the domains in this formula.
A domain of type R with local multiplicity $k$ has $2k+2$ moving
coordinates and decomposes into $2k+1$ empty rectangles.  One obtains
this decomposition by starting at the initial state's point on the
vertical circle through the distinguished intersection $p$, following
the first horizontal boundary segment, and removing the rectangle
against that circle.  Repeat with the new initial point on that same
circle.  The local zero-multiplicity condition forces each rectangle to
be empty and determines it uniquely.  Two successive removals reduce
the three equal multiplicities at $p$ by one and remove two moving
coordinates.  The last step is a rectangle, the case $k=0$.  This is
an induction on $k$.  The condition of exactly $2k+2$ moving coordinates
excludes an additional boundary component: following the component
through $p$ has already used all those coordinates.  The same argument
with odd complexity proves the corresponding statement for L domains.
It also proves finiteness: $2k+2$ is at most the grid size, each boundary
segment has finitely many endpoints and two possible circular arcs,
and its local multiplicity at $p$ determines the remaining additive
constant of the domain.

\subsection{The domain cancellation with separate variables}

Here are the cuts proving \eqref{eq:R-full-variable-identity}.  They
are included because retaining the old variables is essential.  Lift
a rectangle of $G$ to the corresponding rectangle of $A$.  Such a
rectangle either avoids $p$, or has $p$ as its sole interior state
point; in the latter case its coefficient still omits the new O.
Compare an R count followed by one of these lifted rectangles with
an empty rectangle of $A$ followed by an R count.  A \emph{moving corner} is a
point in the symmetric difference of the two states at a step.  Let
$m=2k+2$ be the complexity of the R domain.

For clarity the complete cutting list is the following.  In this
list $A_j$ means that the rectangle avoids $p$ and shares $j$ moving
corners with the R domain; $B$ means that $p$ is a corner of the
rectangle; $C$ that $p$ lies on an open rectangle edge; and $D_j$ that
$p$ is the rectangle's sole interior state point and there are $j$
shared moving corners.  A subscript $m$ records the R complexity.
The symbol $C^-$ denotes the rectangle on the left of $p$, and
$C^+$ the rectangle above $p$.  For $C^-$ let $\delta$ be the
innermost horizontal width of the R domain; for $C^+$ use its
innermost vertical height.
\[
\begin{array}{c|c|c}
\text{first presentation}&\text{other presentation}&\text{operation}\\\hline
A_0&A_0&\text{reverse the order of the pieces}\\
A_1&A_1&\text{the other cut at the concave corner}\\
(A_2)_m&(D_0)_{m-2}&\text{transfer two consecutive turns},\quad m\ge4\\
B_4\ (m=2)&\text{the other thin annulus}
       &\text{new-column / old-root-row annuli}\\
(C^-)_m, \delta=1&(D_1)_m&\text{add the new-column annulus}\\
(C^-)_m, \delta>1&(A_3)_{m+2}&\text{add the new-column annulus}\\
(C^+)_m, \delta=1&(D_1)_m&\text{add the old-root-row annulus}\\
(C^+)_m, \delta>1&(A_3)_{m+2}&\text{add the old-root-row annulus}
\end{array}
\]
The last two rows also have the L-domain presentation described below;
the $B_4$ row has the trivial L-domain presentation.

We explain both exhaustiveness and the indicated constructions.
Away from $p$, three zero quadrants at every state point permit zero,
one, or two shared moving corners.  Zero corners allow the two pieces
to be removed in either order.  One shared corner produces a unique
concave corner, with the two cuts indicated in the table.  With two
shared corners, following the boundary removes two successive turns
of the R domain and puts $p$ into the lifted rectangle; reversing
that cut restores those two turns.  This is exactly the $A_2/D_0$
pair.  Three shared corners force the fourth side to go once around
a row or column circle: the rectangle has width or height one and
the associated innermost R width or height is one.  These are the
$A_3$ entries, recovered uniquely by deleting the specified annulus.

If the extra rectangle has $p$ as a corner, its two intermediate
coordinates are the unique state points on the two circles through
$p$.  The R zero-quadrant conditions force $m=2$, with the same
two moving coordinates.  The composite is consequently a thin
annulus, either the new-column annulus or the row annulus containing
$O_0$.  If $p$ is instead on an open edge, the same quadrant conditions
allow precisely a rectangle to the left or above $p$; the bottom and
right alternatives would leave an initial state point in a positive
quadrant of the R domain.  These are $C^-,C^+$.  Adding the annulus
shown in the table completes the boundary across $p$.  When
$\delta=1$, the completed boundary cuts off a rectangle with $p$ as
its sole interior state point.  When $\delta>1$, it cuts off a thin
rectangle avoiding $p$ and introduces two further turns in the R
boundary.  The cut starts at the initial state's point in the
inserted row or column, so it is unique; deleting that same annulus
is the inverse construction.  A lifted rectangle having $p$ inside
can share only zero or one moving corner, giving exactly $D_0,D_1$.
Finally an empty input rectangle with $p$ strictly inside cannot
precede an R domain: follow the R boundary from its initial point on
the vertical circle through $p$ to the first exit from that rectangle.
Either this initial point is already inside the rectangle, violating
emptiness, or that first exit puts one of its corner state points in
a positive quadrant of the R domain, violating the three-zero rule.
Thus the table accounts for every contribution.

For the first three rows the complete two-chain is unchanged by the
cut, and so is every old marking multiplicity.  The same holds for
the $C^-$ rows except for adding the new-column annulus.  That
annulus contains only $O_i$, whose multiplicity in the R factor is
omitted from the coefficient.  Hence these contributions cancel in
pairs with their full old-variable monomials intact.

For $B_4$, the two annuli have coefficients $U_i$ and $U_0$.  They
give $(U_i+U_0)$ times the trivial L-domain coefficient.  In a $C^+$
row the underlying composite itself has three equal multiplicities
$k+1$ at $p$ and southwest multiplicity $k$; it is an L domain of
complexity $m+1$.  Let $w$ be its monomial omitting $U_i$.  The
rectangle in the $C^+$ presentation contains $O_i$, so its
contribution is $U_iw$.  Adding the old-root-row annulus gives the
other presentation, with contribution $U_0w$.  Their sum is exactly
$(U_i+U_0)w$.  Every nontrivial L domain is obtained once in this
way: cut at its innermost horizontal boundary segment through the
distinguished circle, leaving the even-complexity R domain.  This
is the inverse of the $C^+$ construction.  Together with the trivial
case it accounts for the whole right side of
\eqref{eq:R-full-variable-identity}.

No argument in this cancellation used X markings.  It needs only
one O per row and column, the new O in the northeast square, and
the designation of the old O in the row immediately below it.
It therefore also holds when that old O happens to occupy the
southwest square.  This observation justifies the temporary O-only
normal form used below, as well as auxiliary coincident markings
elsewhere.  Substituting $U_i=U_0$ proves the claimed multivariable
chain equation without identifying any further variables.

\subsection{Gradings and polynomial coefficients}

Cyclically choose coordinates so that $G$ has size $n$, the stabilized
X lies in its last column and row $0$, and the new column and new row
of $A$ are column $n$ and row $1$.  Thus $p=(n,1)$.  If $y$ contains
$p$ and $e(y)$ deletes that coordinate, direct substitution in $J$ gives
\[
 M_A(y)=M_G(e(y))-1,\qquad A_A(y)=A_G(e(y))-1.
\]
For the first equality, the extra contributions to
$J(y,y),-2J(y,O_A),J(O_A,O_A)$ are respectively $1,-3,1$.
The analogous change of $M_X$ is zero, proving the second equality
after the change of grid size is included.  These calculations use
only the row and column permutations.  They remain true if an old
O coincides with the chosen old X.

Let $q:x\to y$ be a type-R domain with three local multiplicities $k$.
Its multiplicity at the new O is $k$, and that at the southeast X is
$k+1$.  Put $w(q)=\sum_{j\ne i}n_{O_j}(q)$.  Summing the rectangle
grading formulas over the decomposition above gives
\[
 M_A(x)-M_A(y)=2k+1-2(k+w(q))=1-2w(q).
\]
Consequently $T^{w(q)}e(y)$ has exactly the degree of $x$.  For the
second filtration coordinate the same calculation gives
\[
 A_G(e(y))-w(q)-A_A(x)
   =1+k-n_X(q)
   =-\!\sum_{X\ne X_{\rm SE}}n_X(q)\le0.
\]
All exponents $w(q)$ are nonnegative.  Thus $R$ has degree zero and
shifts $(0,0)$, also for auxiliary marking pairs with other coincident
O and X markings.

\subsection{Compatibility with O switches away from the new column}

The following argument concerns only O diagrams, so a temporary
coincidence with an X marking has no effect.  Exchange O markings in
two adjacent old columns.  They are still adjacent in $A$, since its
new column is last.  Let $F_G,F_A$ be the respective diagonal O-switch
maps, and let $R,R'$ be the type-R counts before and after this change.
Their domain sets are identical: the definition of a type-R domain
does not refer to O markings.

Write $\varepsilon_A,\varepsilon_G$ for the switch indicators on the
two grids, reserving $e(y)$ for deletion of $p$.  For any such domain
$q:x\to y$, the same boundary computation as for a rectangle gives
\[
 n_{O'}(q)-n_O(q)=\varepsilon_A(y)-\varepsilon_A(x).
\]
It can equivalently be obtained by summing the rectangle identities
over its decomposition.  The new O has not moved.  Since the switch
uses old columns, $\varepsilon_A(y)=\varepsilon_G(e(y))$.  Comparing
the monomial of each individual domain therefore gives the exact identity
\begin{equation}\label{eq:R-switch-exact}
 R'F_A=F_G R.
\end{equation}
The same assertion holds for the reverse diagonal switch map.

Choose, among the two exchanged labels, a label $j\ne0$.  There is
always such a label.  Differentiation of
\eqref{eq:R-full-variable-identity} after $U_i=U_0$ proves that $R$
and $R'$ intertwine $\Phi_j$ on homology.  Using
\eqref{eq:R-switch-exact} in the definition of the degree-zero
Clifford maps $K_j$ consequently gives
\begin{equation}\label{eq:R-K-naturality}
 R'_* (K_j^A)_*=(K_j^G)_* R_*.
\end{equation}
This remains valid when label $0$ itself is one of the exchanged
markings: the detecting label $j$ is the other one, so its variable
remains independent of $U_i=U_0$.

\subsection{An explicit normal-form evaluation}

Bubble-sort the old O permutation to $O_G(c)=n-1-c$, using adjacent
old columns only.  Each $K_j$ in \eqref{eq:R-K-naturality} preserves
the canonical class and is invertible on homology.  It follows that
the desired nonvanishing for $R$ is equivalent to nonvanishing in
this one normal form.  The old O in row $0$ is then in the last
column; hence the chosen old X would coincide with it.  This temporary
O-only calculation is legitimate: the type-R domain identity and
\eqref{eq:R-K-naturality} use no X markings, and canonical classes
are defined from the O complex using any disjoint X permutation.

The stabilized O permutation in this normal form is
\[
 O_A=(n,n-1,\ldots,2,0,1).
\]
It is obtained from the size-$(n+1)$ antidiagonal by interchanging
its last two O markings.  Write
\[
 \theta=(0,n,n-1,\ldots,2,1),\qquad
 \theta_G=(0,n-1,n-2,\ldots,1).
\]
The antidiagonal canonical cycle is $\theta$.  For this last O switch
the degree-minus-one diagonal map is
$F(x)=T^{1-\mathbf1_{\{x(n)=1\}}}x$, so it sends $\theta$ to itself.
Every derivative of $\theta$ in the antidiagonal complex is zero
already on chains, since it would have positive degree and polynomial
coefficients, while that complex has no positive-degree polynomial
generators.  The Clifford identity therefore identifies the canonical
cycle of $A$ as
\[
 z=\Phi_i^A\theta.
\]

Here this cycle can be written completely explicitly.  For
$0\le b\le n-1$, let $x_b$ be obtained from $\theta$ by interchanging
its coordinates in columns $b,n$.  The empty rectangles from
$\theta$ through the new O are exactly the $n$ rectangles with
southwest corner $(n,1)$ and northeast corner $(b,\theta(b))$,
using the circular intervals from $n$ to $b$ and from $1$ to
$\theta(b)$.  Indeed, a rectangle through that O which did not start
on column $n$ would contain the state point $(n,1)$ in its interior.
For the displayed rectangles all intervening antidiagonal state
points and all old O markings lie outside.  Each contains the new
O once and no other O.  Hence
\[
 z=\sum_{b=0}^{n-1}x_b.
\]
Every $x_b$ has Maslov degree zero.  The only possible degree-zero
polynomial output of $R$ in the old antidiagonal complex is
$\theta_G$ with coefficient $1$, since its other states have strictly
negative Maslov degree.  The corresponding stabilized terminal state
is precisely $\theta$.

Since $x_b$ differs from $\theta$ in two coordinates, a type-R domain
from $x_b$ to $\theta$ must have $k=0$; it is a rectangle with
northwest corner $p=(n,1)$.  For $b=0$ the unique such rectangle is
the single square in the new column and row $0$.  It contains no O
marking and contributes $\theta_G$.  For every $b\ge1$, the only
rectangle with the required southeast multiplicity at $p$ uses the
circular intervals from column $n$ to column $b$ and from row
$\theta(b)$ to row $1$.  It contains the unchanged state point
$(0,0)$ strictly in its interior, so it is not a type-R domain.
Therefore
\[
 Rz=\theta_G.
\]
This proves nonvanishing in the normal form.  Transport through
\eqref{eq:R-K-naturality} proves it for every O permutation.  Since
$R$ has degree zero and preserves the polynomial subcomplex, Lemma N
now yields
\begin{equation}\label{eq:R-canonical-independent}
 R_*\operatorname{can}(A)=\operatorname{can}(G).
\end{equation}
This establishes the property of this specific polynomial formula
without using any unspecified stabilization map or any comparison
with holomorphic Floer homology.

\subsection{The outgoing top map and colour naturality in both directions}

For completeness, the stabilization direction has an equally explicit
polynomial representative.  Given a state $x$ of $G$, extend it by $p$
to a state $\widetilde x$ of $A$.  Define
\[
 Sx=\sum_y\sum_{q:\widetilde x\to y}
       \prod_{j\ne i}U_j^{n_{O_j}(q)}y,
\]
where the inner sum is over nonnegative outgoing domains satisfying
the same three-zero rule away from $p$, having multiplicity $k+1$
in the northeast square at $p$ and $k$ in the other three squares,
and having $2k+2$ moving coordinates.  These domains have the same
finite boundary description and rectangle decomposition as the
incoming domains, with arrows reversed and the local vertical
coordinate reflected.  This description is a definition of the
formula, not an instruction to reflect the grid or its complex.

The boundary cuts in the preceding table apply with their arrows
reversed and with north and south interchanged.  The two annuli in
the exceptional rows are now the new-column and new-row annuli.
Each contains only the new O, and in each presentation its
multiplicity belongs to the outgoing domain, whose monomial omits
that O.  Their coefficients are therefore equal.  More explicitly,
the reflected $C^+$ configuration has an extra rectangle below $p$.
That rectangle contains no new O.  Its union with the outgoing
domain has northeast, southeast, and southwest multiplicities
$k+1$ and northwest multiplicity $k$.  Adding the new-row annulus
changes the northeast multiplicity to $k+2$ and the other three to
$k+1$.  These are the two alternative presentations; both have
exactly the same old-variable monomial.  The complexity-two
exception has the same two equal coefficients.  The other rows
cancel by equal total domains or by the new-column annulus, exactly
as before.  Consequently
\begin{equation}\label{eq:S-full-variable-identity}
 D_A S=S D_G
\end{equation}
already with \emph{all} O variables independent, interpreting the
source as extended by the unused variable $U_i$.

For an outgoing domain with parameter $k$, its multiplicity at
$O_i$ is $k+1$ and at the southeast X is $k$.  Its $2k+1$-rectangle
decomposition and the displayed state-extension grading formulas give
\[
 M_A(y)-2\sum_{j\ne i}n_{O_j}(q)=M_G(x),
\qquad
 A_A(y)-\sum_{j\ne i}n_{O_j}(q)-A_G(x)
       =k-n_X(q)\le0.
\]
Thus $S$ has degree zero and shifts $(0,0)$.

Its canonical-class action follows by the same normal-form method,
and can be checked particularly simply.  The O-switch boundary
identity gives
\[
 F_A S=S'F_G
\]
term by term.  Differentiating \eqref{eq:S-full-variable-identity}
and using this exact identity gives the corresponding $K$-naturality
on homology.  Therefore sort the old O permutation as above and
evaluate $S$ on $\theta_G$.  Its extension is $\theta$.

An outgoing domain starting at $\theta$ cannot have $k>0$.
Indeed, the southeast unit square at $p=(n,1)$ is the northwest
square at the initial state point $(0,0)$, using the circular last
column.  Its multiplicity is $k$.  The three-zero rule at $(0,0)$
would then leave a positive northwest quadrant alone.  This has
the boundary sign of a terminal corner, whereas $(0,0)$ is an
initial state point; if that point were also terminal, all four
quadrants would have to vanish.  Both possibilities contradict
$k>0$.  Hence $k=0$.  The domains are precisely the $n$ rectangles
from $\theta$ to $x_b$ already enumerated above.  Their only O is
the new one, whose factor is omitted.  We conclude
\[
 S\theta_G=\sum_{b=0}^{n-1}x_b=z.
\]
Transport by $K$ proves that this specific map $S$ preserves the
canonical class for every old O diagram.

Finally identify the new O variable with the variable of the
component on which the stabilization is performed, and give every
other component its own variable.  Equations
\eqref{eq:R-full-variable-identity} and
\eqref{eq:S-full-variable-identity} are now chain equations over
that colour ring in both directions.  Differentiating in any colour
variable proves naturality for the sum of its marking derivatives.
For a commutation the empty-pentagon formula retains the complete
O monomial of each polygon; the alternate rectangle--pentagon cut
has that same monomial, so its chain equation is valid over the
colour ring too.  Cyclic permutations simply relabel these
monomials.  The granted corner reduction expresses every other
stabilization orientation using these corner maps and such
commutations and permutations.  Every extra O in that reduction is
on the same physical component as the O with which its variable is
identified.  Composing the differentiated identities therefore proves
colour naturality for every ordinary move used in the construction,
in both directions, on arbitrary homology classes.

\end{proof}


\section{Explicit relations for signed word grids}

\begin{lemma}[Signed word-grid relations; new combinatorial]
\label{lem:signed-relations}
\formalname{signedWordGridRelations}
Every finite derivation by Artin relations, far commutations, free
inverse cancellation or insertion, and cyclic whole-word rotations
has a finite ordinary-grid realization under the signed gadget
convention above.  All local replacements work in arbitrary signed
contexts with the indicated port incidences.  A rank becoming
unused is retained as the specified four-mark block.  For an
Artin/free derivation whose initial and final words contain every
generator, the external guard
in the proof makes every word in the initial normalization contain
every generator.  All displayed rules specify legal noninterleaving
exchanges and genuine three-corner stabilization patterns.
\end{lemma}
\begin{proof}

The negative-letter gadget has outgoing heights
$(r+\tfrac12,r+1,r)$ and incoming heights
$(r+1,r,r+\tfrac12)$ in its three consecutive columns.  Integer-height
arcs join consecutive occurrences cyclically, exactly as for positive
letters.  The following constructions work in an arbitrary signed-word
context, and use no topological equivalence theorem.

\subsection{Two elementary grid rules}

We distinguish strict noninterleaving commutations from the case of a
shared endpoint.  The latter case can also be performed by ordinary grid
moves, as follows.  Suppose consecutive columns $A,B$ have
\[
 (O_A,X_A)=(r,p),\qquad (O_B,X_B)=(q,r).
\]
First assume $p\ne q$.  The row $r$ joins the consecutive columns $A,B$.
Move this row toward row $p$, along the arc of the row circle from $r$
to $p$ which avoids $q$, stopping adjacent to $p$.  Every row actually
crossed has both its markings outside $A,B$, so each exchange is a strict
noninterleaving row commutation.  The block with columns $A,B$ and rows
$r,p$ has precisely the three markings
\[
 O(A,r),\quad X(A,p),\quad X(B,r).
\]
Destabilize, deleting column $A$ and row $r$, and replacing the two $X$'s
by $X(B,p)$.  Now stabilize this $X$ using a new column on the opposite
side of $B$: restore row $r$ and put
\[
 (O_B,X_B)=(q,r),\qquad (O_A,X_A)=(r,p)
\]
in the order $B,A$.  Move the short row $r$ back through the same rows
in reverse order.  This is exactly the required exchange of $A,B$.

If $p=q$, the two columns form a two-row unknot block.  Insert a row $t$
next to $p$ and a column $D$ immediately before $A$, by the $X$
stabilization
\[
 D:(O=t,X=p),\quad A:(O=r,X=t),\quad B:(O=p,X=r).
\]
Exchange $A,B$ by the preceding rule, then exchange $D,B$ by its
$O/X$-interchanged version.  Both exchanges have distinct nonshared
endpoints.  The order is now $B,D,A$.  Destabilize the block
$O(D,t),X(D,p),X(A,t)$, deleting $D,t$.  The result is $B,A$ with exactly
the required markings.  All other possible shared-endpoint patterns
follow by interchanging $O,X$ and/or reversing the names of the two
columns.  Transposing this rule proves its row version.  Henceforth
``exchange'' means a strict commutation or this fully specified
replacement when endpoints coincide.

The second rule contracts a bivalent column.  Suppose column $C$ carries
$O$ in row $u$ and $X$ in a different row $v$, with $u,v$ adjacent, and
let $E$ be the other column carrying an $X$ in row $u$.  Slide $C$
around the column circle until immediately before $E$.  Its marking
interval has consecutive endpoints, so all intervening exchanges are
noninterleaving or shared-endpoint exchanges just treated.  If
$O(E)\ne v$, destabilize the three-corner block
\[
 O(C,u),\qquad X(C,v),\qquad X(E,u),
\]
deleting $C,u$ and replacing $X(E,u)$ by $X(E,v)$.  This joins the two
horizontal arcs.  If $O(E)=v$, the block is instead an isolated $2\times2$
unknot and is retained as such.

The inverse rule subdivides any horizontal arc: at its $X$ endpoint
insert an adjacent row and column using the reverse displayed
three-corner pattern, then slide the resulting short vertical column to
the desired position.  Thus an incoming and an outgoing port which were
the same horizontal arc can always be separated by an ordinary
stabilization, with a bivalent column outside a specified local word
block.  After a local construction this subdivision can be contracted
unless that entire strand has become unused, in which case its
$2\times2$ unknot block is retained.

\subsection{Choice of row orders at a repeated height}

At an integer height the horizontal arcs run between successive
occurrences of that rank.  Their interiors in the column circle are
disjoint.  At a half-integer height the arcs lie in distinct three-column
letter blocks and again have disjoint interiors.  Therefore any two
rows of the same height can be exchanged after bringing them next to
one another: their endpoint pairs are nonalternating, with shared
endpoints covered by the first rule above.  Bubble-sort consequently
realizes every desired tie-break by ordinary grid moves.

\subsection{Cancellation of consecutive inverse letters}

Consider $\sigma_r\sigma_r^{-1}$ and name its six consecutive columns
$A_0,A_1,A_2,B_0,B_1,B_2$.  Denote its horizontal rows by
\[
\begin{array}{c|c|c}
\text{height}&\text{row}&\text{arc}\\\hline
r&a&\text{previous occurrence}\to A_0\\
r&b&A_1\to B_1\\
r&c&B_2\to\text{next occurrence}\\
r+\tfrac12&p&A_0\to A_2\\
r+\tfrac12&q&B_0\to B_2\\
r+1&d&\text{previous occurrence}\to A_1\\
r+1&e&A_2\to B_0\\
r+1&f&B_1\to\text{next occurrence}.
\end{array}
\]
Thus the marking lists in these six columns are
\[
 O=(p,b,e,q,f,c),\qquad X=(a,d,p,e,b,q).
\]
If $a=c$ or $d=f$, subdivide that horizontal arc outside this six-column
block by the second elementary rule.  These rows are then distinct.
The subdivision is remembered as belonging to a rank which will become
unused after the cancellation.

Choose $p,q$ to be the last two rows at height $r+\tfrac12$, in that
order, and $e$ to be the first row at height $r+1$.  Perform:
\begin{enumerate}
\item The rows $q,e$ and columns $A_2,B_0$ contain
      $O(A_2,e),X(B_0,e),O(B_0,q)$.  Delete $B_0,e$ by
      destabilization, replacing these by $O(A_2,q)$.
\item Exchange $A_2,B_1$.  The former has endpoints $p,q$, while the
      latter has endpoints $b,f$, so the intervals are nested.  Now
      $A_1,B_1$ are adjacent and row $b$ is a short row.  Move it next
      to $d$.  Destabilize
      $O(A_1,b),X(A_1,d),X(B_1,b)$, deleting $A_1,b$ and putting
      $X(B_1,d)$.
\item Destabilize the block
      $O(A_2,q),X(A_2,p),X(B_2,q)$, deleting $A_2,q$ and putting
      $X(B_2,p)$.  The rows $p,q$ are still adjacent.
\item Exchange $B_1,B_2$.  Both endpoints of $B_1$ are in the upper
      integer-height group, and the interval of $B_2$ runs from the
      lower group to $p$, so they are disjoint.  Now row $p$ has the
      adjacent endpoints $A_0,B_2$.  Move it next to $c$ and
      destabilize
      $O(A_0,p),X(B_2,p),O(B_2,c)$, deleting $B_2,p$ and putting
      $O(A_0,c)$.
\item The remaining two local columns are
      $A_0:(O=c,X=a)$ and $B_1:(O=f,X=d)$.  Bring $a,c$ together
      within their common height group and apply bivalent contraction
      to $A_0$; do the same with $d,f$ and $B_1$.
\end{enumerate}
For a rank which still occurs outside the cancelled pair, the final
contraction is a three-corner destabilization.  For an unused rank, the
remembered subdivision column and this final column form precisely an
isolated $2\times2$ unknot; retain that block as the resolution of the
unused strand.  There is no death or saddle in this case.  The two
columns of this block can be brought together, and the block can be
moved to any prescribed auxiliary position using the elementary
exchanges above.

When both ranks persist, the construction uses exactly six
destabilizations, with exchanges and harmless tie-break changes.  All
six gadget columns have disappeared, and the two surviving
integer-height arcs are exactly the arcs formed by skipping the two
cancelled occurrences.  Columns outside the block retain their cyclic
order; if necessary make one cyclic column permutation to restore the
chosen cut.  The resulting grid is the signed word grid with the pair
removed, up to tie-break changes.  This checks the endpoint by marking
incidence, not by a link-isotopy assertion.

For $\sigma_r^{-1}\sigma_r$, reverse the local row coordinate, applying
the same rule with upper and lower ranks interchanged.  This means
conjugating each explicit move by row reflection, not performing a
reflection as a grid move.  Commutations and the eight stabilization
patterns are closed under that reflection, so the reflected sequence is
again a sequence of the permitted moves.  Inverse-pair insertion is the
reverse of the corresponding sequence.

\subsection{The three-letter Artin relation}

For the source block $\sigma_r\sigma_{r+1}\sigma_r$ label its columns
$0,1,\ldots,8$.  Each outgoing row is labelled by its starting column.
Let $9,10,11$ denote the incoming port rows at integer heights
$r,r+1,r+2$.  If any such port is the same arc as its outgoing port,
first subdivide that arc outside the block as above, so all twelve
distinguished rows are distinct.  No other subdivision is required.

The initial $X$ incidences in the nine columns are
\[
 X=(9,10,0,2,11,3,1,4,6),\qquad O(c)=c.
\]
Choose tie-breaks so that the distinguished rows have the relative order
\[
 (1,7,9,\ 0,6,\ 2,4,8,10,\ 3,\ 5,11).
\]
The groups displayed correspond respectively to heights
$r,r+\tfrac12,r+1,r+\tfrac32,r+2$.  Outgoing port rows $7,8,5$ end at
columns outside the block; incoming port rows $9,10,11$ start at
columns outside the block.

In the following table $C_{a,b}$ exchanges the columns with persistent
labels $a,b$, and $R_{a,b}$ exchanges the rows with those labels.  Each
column pair is adjacent at its indicated step.  For rows, first slide
the internal row across any nondistinguished rows between the indicated
pair, then perform the exchange.  Such intermediate rows have both
endpoints outside the local column block, so these preliminary slides
are strict noninterleaving commutations.

The last two columns verify noninterleaving at every step.  For a column
exchange the numbers are positions in the current order of the twelve
distinguished rows.  For a row exchange they are positions in the
current order of the nine local columns; $*$ is any outside column, all
of which lie in the complementary interval after position $8$ when
the circle is cut immediately before the local block.  The endpoint
pairs are disjoint, nested, or share an endpoint.  Shared endpoints use
the fully specified elementary rule above.
\[
\begin{array}{r|c|c|c}
1&C_{0,1}&(2,3)&(0,8)\\
2&R_{1,7}&(0,6)&(7,*)\\
3&R_{6,2}&(6,8)&(2,3)\\
4&C_{5,6}&(9,10)&(1,5)\\
5&C_{5,7}&(9,10)&(0,6)\\
6&C_{5,8}&(9,10)&(5,7)\\
7&C_{2,3}&(3,4)&(4,9)\\
8&C_{2,4}&(3,4)&(6,11)\\
9&C_{2,6}&(3,4)&(1,5)\\
10&C_{2,7}&(3,4)&(0,6)\\
11&C_{4,6}&(6,11)&(1,5)\\
12&R_{6,4}&(3,7)&(4,5)\\
13&R_{2,4}&(2,6)&(4,5)\\
14&R_{6,8}&(3,7)&(7,*)\\
15&R_{6,10}&(3,7)&(0,*)\\
16&R_{6,3}&(3,7)&(2,8)\\
17&C_{3,6}&(5,8)&(1,9)\\
18&C_{0,6}&(2,3)&(1,9)\\
19&R_{1,9}&(0,1)&(2,*)\\
20&R_{1,0}&(0,1)&(2,6)\\
21&C_{3,4}&(5,8)&(4,11)\\
22&C_{0,4}&(1,2)&(4,11)\\
23&R_{1,4}&(0,1)&(2,5)\\
24&C_{6,4}&(4,9)&(3,11)\\
25&C_{3,7}&(5,8)&(0,3)\\
26&R_{1,2}&(0,2)&(5,6)\\
27&C_{3,2}&(4,8)&(2,4)\\
28&R_{1,8}&(0,2)&(7,*)\\
29&R_{1,10}&(0,2)&(0,*)
\end{array}
\]
Every row exchange includes an internal row: these are rows
$0,1,2,3,4,6$, whose two endpoints always remain among the nine local
columns.  Thus the preliminary slides past entirely external rows are
always permitted.  In rows of the table involving a port, the displayed
noninterleaving condition holds for every possible outside endpoint,
not just for one choice of closure.

After step 29 the column order is
\[
 (1,4,6,0,7,2,3,8,5)
\]
and the distinguished row order is
\[
 (7,9,\ 0,\ 4,2,8,10,\ 1,3,\ 6,5,11).
\]
The five groups now have heights
$r,r+\tfrac12,r+1,r+\tfrac32,r+2$.  Reinsert the entirely external rows
in height order: only internal rows need to pass them, and those
exchanges are harmless as just proved.  The three incoming and three
outgoing port rows never change their heights or their order relative
to external rows.

To check the endpoint, identify the new three-letter columns with the
old labels by
\[
 (\text{new }0,\ldots,8)\longmapsto(1,4,6,0,7,2,3,8,5).
\]
The unchanged $O/X$ incidence, with the final row-height groups just
listed, is precisely the gadget sequence
$\sigma_{r+1}\sigma_r\sigma_{r+1}$.  In particular the three incoming
ports and the three outgoing ports are the same, so this is a
replacement in an arbitrary surrounding word.  Contract any temporary
port subdivisions, then choose the specified final tie-break.  This
finishes a completely local grid proof of the Artin relation.

The all-negative relation follows by reflecting the local height
coordinate and conjugating each listed grid move by that reflection.
Alternatively it follows by inserting inverse pairs, applying the
positive relation, and cancelling inverse pairs.  No mixed-sign
relation beyond these generators of the group congruence is needed.

\subsection{Far commutation and cyclic rotation}

If $|i-j|\ge2$, every vertical marking interval in the three columns
of a $\sigma_i^{\pm1}$ gadget lies in the height strip $[i,i+1]$, and
those of a $\sigma_j^{\pm1}$ gadget lie in $[j,j+1]$.  These strips are
disjoint.  Move the three columns of the second block past the three
columns of the first block, one at a time, for nine strict column
commutations.  The rank sets are disjoint, so the integer-arc incidence
does not change except for the indicated block order.  This proves far
commutation for every choice of signs.

For a cyclic word rotation $uv\mapsto vu$, cyclically permute the
columns by $3\ell(u)$.  Integer-arc incidence is cyclic and hence
unchanged, and the row-height order is unchanged.  Correcting the
tie-break by row exchanges gives the chosen word grid.  Any auxiliary
unused-strand blocks can be kept at their prescribed positions by
sliding their short rows and columns as above.

Consequently every finite derivation using free inverse cancellation,
Artin relations, far commutation, and cyclic word rotations has an
explicit finite realization by ordinary grid moves.  Each generator
of the derivation is replaced by the rule specified here; there is no
appeal to the topological statement that isotopic link grids are
connected by moves.

\subsection{Moving a consecutive active block to the first ranks}

Set $\rho=\sigma_0\sigma_1\cdots\sigma_{S-2}$.  For
$0\le i\le S-3$ the following calculation uses only the two positive
Artin relations just realized:
\begin{align*}
 \rho\sigma_i
 &= (\sigma_0\cdots\sigma_{i-1})
       \sigma_i\sigma_{i+1}\sigma_i
       (\sigma_{i+2}\cdots\sigma_{S-2})\\
 &= (\sigma_0\cdots\sigma_{i-1})
       \sigma_{i+1}\sigma_i\sigma_{i+1}
       (\sigma_{i+2}\cdots\sigma_{S-2})\\
 &=\sigma_{i+1}\rho.
\end{align*}
The first equality moves the final $\sigma_i$ left past indices at
least $i+2$; the last moves $\sigma_{i+1}$ left past indices at most
$i-1$.  Empty strings at the ends are omitted.  Appending
$\rho^{-1}$ and freely cancelling proves
$\rho\sigma_i\rho^{-1}=\sigma_{i+1}$ by a finite word derivation.
Taking inverses gives the same identity for $\sigma_i^{-1}$.
Consequently, if $\beta$ uses only indices $0,\ldots,a-2$ and
$a+k\le S$, iteration gives
\[
 \operatorname{shift}_k(\beta)=\rho^k\beta\rho^{-k}.
\]
This identity is obtained letter by letter, inserting inverse pairs
between consecutive letters when needed.  It therefore has an
explicit ordinary-grid realization by the preceding rules.  In
particular it applies to a one-crossing factor, as well as to a local
positive word on any consecutive block of active strands.

\subsection{A guard preserving the ambient number of strands}

When a finite word derivation is to be realized before the scissor
construction, its intermediate grids can all be kept on the prescribed
$S$ ranks, with every generator present, as follows.  Put
$p=\rho\rho^{-1}$.  This literal word contains every generator and
hence uses every rank, and represents the identity.  Suppose $\alpha$
contains every generator and a finite derivation by local Artin and free
relations changes $\alpha$ into a product
$Q=q_1\cdots q_r$, where every $q_j$ has been padded as
$p\widetilde q_jp$.  Perform the derivation in the sequence
\[
 \alpha\longrightarrow p\alpha
 \longrightarrow pQ\longrightarrow Q.
\]
The first arrow inserts inverse pairs in front of $\alpha$; its
original letters keep every generator present during that insertion.  In
the middle arrow, keep the displayed external copy of $p$ unchanged
and perform the chosen derivation entirely in the remaining suffix.
The local grid rules allow arbitrary surrounding words, so this is
exactly the prescribed realization with a fixed guard.  It is not
necessary to rotate this external guard through any factor.  For the
last arrow cancel its central inverse pairs successively.  The
individually padded suffix $Q$ keeps every generator present until that
external copy has disappeared.

If a preceding comparison is only up to whole-word conjugacy, first
absorb that conjugacy into every factor conjugator:
$c(q_1\cdots q_r)c^{-1}=\prod_j(cq_jc^{-1})$.
Inserting $c^{-1}c$ between factors is a finite free derivation.
This gives exactly the Artin/free suffix derivation used here; no
cyclic rotation of a proper substring is being asserted to be a
word relation.
The conjugator can itself be read from any finite comparison that
uses whole-word rotations.  Maintain an identity
$\alpha=cvc^{-1}$, starting with $c$ empty and $v=\alpha$.
A local Artin or free step changes $v$ and leaves $c$ fixed.
For a rotation $v=ab\mapsto ba$, replace $c$ by $ca$:
\[
 c(ab)c^{-1}=(ca)(ba)(ca)^{-1},
\]
where the equality is realized by inserting $aa^{-1}$ immediately
after the displayed $b$.  This records a finite derivation using
only local Artin and free relations.  At the end absorb the recorded
$c$ into the factor conjugators as above, before inserting the guard.

Thus the guard is absent before any scissor, factor exposure, or marked
conjugacy cleanup.  Each remaining factor has the form
\[
 p\gamma\eta\gamma^{-1}p
   =(p\gamma)\eta(p\gamma)^{-1},
\]
as a literal word, because $p^{-1}=p$.  Its padding is consequently
removed by the same marked inverse-pair cancellation as its
conjugator.  No extra cyclic rotation of a guard, or assumption about
a marked component during the guarded derivation, is required.

\end{proof}


\section{Cancelling inverse letters while retaining an attached summand}

\begin{lemma}[Marked inverse cancellation; new combinatorial]
\label{lem:marked-cancellation}
\formalname{markedInverseCancellation}
Let a closed signed word have a marked horizontal integer-rank arc
in a nontrivial permutation cycle, with an arbitrary prepared grid
$H$ inserted in that arc in the frame specified below.  A consecutive
inverse pair at the marked cut can be cancelled by the displayed
ordinary moves while preserving all internal marking data of $H$.
The mark changes by the cancelled letter's transposition, and the
exact frame is restored.  Iteration reduces a marked
$\gamma\eta\gamma^{-1}$ to $\eta$, with tracked rank
$\pi(\gamma)^{-1}(k)$ from initial rank $k$.  Unused unmarked ranks
become precisely the stated four-mark components.  These moves
are natural for component-colour derivatives.
\end{lemma}
\begin{proof}

Here a protected grid $H$ may contain other components in addition to its
marked component.  Use its supplied seam preparation, so that
\[
 X_H(0)=0,\qquad O_H(h-1)=0.
\]
These are exactly the right-hand preparation conditions for the binary
seam.  No further reordering of the protected grid is required.  In
particular, its last X need not be in its top row.

\subsection{The frame and its elementary exchanges}

In a base grid choose a horizontal arc from $O(P,b)$ to $X(Z,b)$.
The framed insertion of $H$ into that arc has the following data:
the base column $Z$ is retained, with its X moved from $b$ to a new row
$t$; the first column of $H$ has X in row $b$; its last column has O
in row $t$; and every other marking of $H$ is in an internal H row.
Thus the two exterior rows are
\[
 b:P\longrightarrow H_{\rm first},\qquad
 t:H_{\rm last}\longrightarrow Z.
\]
The row circle is ordered as
\begin{equation}\label{eq:exact-marked-frame}
 (t,\ \hbox{all other base rows in their old cyclic order after }b,
       \ b,\ H_1,\ldots,H_{h-1}).
\end{equation}
The H columns form one contiguous block and retain their internal order.

This frame is obtained from an existing seam by explicit ordinary moves.
The marked column $Z$ of the base grid has been removed in that seam:
$X(Z,b)$ is represented by the first H column's X, and $O(Z,u)$ by the
last H column's O.  O-stabilize at this last O, inserting a column $C$
immediately to its right and a row $t$ next to $u$, with
\[
 O(H_{\rm last},t),\quad X(C,t),\quad O(C,u).
\]
Regard $C$ as the restored base column $Z$.  Its O is exactly the O that
the base column originally had.  The row $t$ has consecutive column
endpoints $H_{\rm last},C$, so it can be moved to immediately after $b$
in the cyclic order of the base rows, using strict exchanges and the
explicit shared-endpoint rule.  Then gather all H internal rows between
$b,t$, preserving their order.  To justify every exchange in this last
operation, an H internal row has both endpoints in the H column block.
An exterior row has both endpoints outside that block, or one endpoint
at its first or last column.  These pairs are always disjoint, nested,
or share an endpoint.  Thus every exchange is one already explicitly
realized by ordinary grid moves.  This proves
\eqref{eq:exact-marked-frame} from the seam, without an isotopy theorem.

The frame has an equally explicit inverse.  In an exact frame move the
H block immediately before the restored base column $Z$, using the
column exchanges established below.  The output row $t$ is then short,
with endpoints $H_{\rm last},Z$.  Move it next to $u=O(Z)$ and
O-destabilize the block
\[
 O(H_{\rm last},t),\quad X(Z,t),\quad O(Z,u),
\]
deleting $Z,t$ and putting $O(H_{\rm last},u)$.  This restores exactly
the seam marking formula.  Choose the column cut at the removed marked
column and the base row cut after the input row, then gather the
internal H rows above the base rows.  The same permitted row exchanges
give the actual binary seam, with its prescribed internal row orders.

Two elementary facts will be used repeatedly.
\begin{enumerate}
\item In the exact frame, every H column can be exchanged with every
base column.  Cut rows at $t$ in \eqref{eq:exact-marked-frame}.  All base
column endpoints lie between $t$ and $b$.  Every H column except its
last has both endpoints between $b$ and $H_{h-1}$; the last H column
joins $t$ to its own internal X row, which is above $b$, and therefore
encloses the whole base interval.  The only
endpoint equalities are handled by the shared-endpoint rule.  Hence
the complete H block can be moved to any gap in the base column order,
without changing that order or the order within H.
\item Even after the two exterior port rows have changed, every
internal H row can be exchanged with every exterior row.  The column
argument in the preceding paragraph still applies: the only exterior
endpoints in the H block are at its first and last columns.  Thus the
H internal rows can be placed in any desired interval of the row
circle, preserving the order of all exterior rows and their own
internal order.  If the two port rows are adjacent in the base order,
putting the H internal rows between the input and output ports restores
the exact frame.
\end{enumerate}

The following boundary absorptions also use only three-corner
destabilizations.  If a base column $L$ immediately precedes H and has
$O(L,b),X(L,d)$, while the H input is $X(H_{\rm first},b)$, move the
short row $b$ next to $d$ and delete $L,b$.  The H input becomes
$X(H_{\rm first},d)$.  Dually, if a base column $R$ immediately follows
H and has $X(R,t),O(R,f)$, while $O(H_{\rm last},t)$ is the H output,
move the short row $t$ next to $f$ and delete $R,t$.  The H output
becomes $O(H_{\rm last},f)$.  The other corner is empty in each case.
These are the usual X- and O-destabilization patterns, respectively.

\subsection{Marked lower connector of a positive--negative pair}

Use the six-column notation of the unmarked cancellation theorem:
\[
 O=(p,b,e,q,f,c),\qquad X=(a,d,p,e,b,q)
\]
in columns $A_0,A_1,A_2,B_0,B_1,B_2$.  Insert framed H into the lower
connector $b:A_1\to B_1$.  Thus $B_1$ now has X in the frame's output
row $t$, and the H ports are $b,t$.  Apply the preliminary port
subdivision of the unmarked rule if the unmarked lower ports $a,c$
coincide.  The marked upper ports $d,f$ are distinct in the application:
otherwise the tracked marked component would be the isolated rank
which becomes unused, contrary to its nontrivial permutation cycle.

First do the ordinary destabilization deleting $B_0,e$ and putting
$O(A_2,q)$, and exchange $A_2,B_1$.  Beforehand H may be moved to a
gap outside these columns by the first framing fact.  Move its block
into the gap between $A_1,B_1$.  Absorb $A_1,b$ into its left boundary,
giving H input $d$, and absorb $B_1,t$ into its right boundary, giving
H output $f$.  Thus H now replaces the upper bivalent column which
occurs after the corresponding two stages of the unmarked cancellation.

Place all H internal rows just above the upper integer-height rows
$d,f$, using the second framing fact.  Consequently every endpoint of
an H column is above the lower and middle rows of the remaining
$A_0,A_2,B_2$ portion.  Delete $A_2,q$ by the ordinary three-corner
destabilization, placing $X(B_2,p)$.  Exchange the entire H column block
past $B_2$: its intervals and that of $B_2$ are disjoint.  Row $p$ now
has consecutive endpoints $A_0,B_2$, so perform the ordinary
destabilization deleting $B_2,p$ and putting $O(A_0,c)$.  Finally
contract the unmarked lower bivalent column $A_0$.

If the unmarked rank persists, all six original letter columns have
been removed by six destabilizations.  The H columns and internal
rows are unchanged in either case.  Its
new input and output rows are precisely $d,f$, the two portions of the
upper-rank arc obtained by deleting the inverse pair.  This is the
same insertion on that arc, before the reframing specified below.
If the unmarked lower rank becomes unused, the last bivalent
contraction instead retains $A_0$ and its preliminary subdivision
column as the $2\times2$ auxiliary block of the unmarked rule; the
other five letter columns have been removed.

\subsection{Marked upper connector of a positive--negative pair}

Now insert framed H into $e:A_2\to B_0$, so $B_0$ has X in row $t$.
Subdivide the unmarked upper arc outside the block if $d=f$.
The marked lower ports $a,c$ are distinct by the same nontrivial-cycle
argument.
Move H between $A_2,B_0$.  Absorb $A_2,e$ at its left boundary, giving
H input $p$, and absorb $B_0,t$ at its right boundary, giving H output
$q$.  H now replaces the middle bivalent column occurring after the
first stage of the ordinary cancellation.

Place its internal rows immediately above the middle row $q$ and below
\emph{every} upper integer-height row.  The interval of $B_1$, whose
endpoints are $b,f$, then contains every H column interval.  Exchange
the whole H block past $B_1$.  The columns $A_1,B_1$ are now adjacent,
so the ordinary short-row move and destabilization delete $A_1,b$ and
put $X(B_1,d)$.

Absorb $B_2,q$ at the right boundary of H, giving output $c$.
Its input is still $p$.  All H column endpoints are now below the
upper-rank endpoints $d,f$ of $B_1$, so exchange $B_1$ past the entire
H block by disjoint-interval commutations.  Absorb $A_0,p$ at the left
boundary of H, giving input $a$.  Contract the unmarked upper bivalent
column $B_1$.

Again exactly six original letter columns have disappeared if the
unmarked rank persists.  H is
unchanged internally and is now inserted on the lower-rank arc with
ports $a,c$.  An unused unmarked upper rank is treated by retaining its
ordinary $2\times2$ auxiliary block, consisting of $B_1$ and the
preliminary subdivision column; the other five letter columns have
disappeared.

\subsection{Reframing and iteration}

The port rows $d,f$ in the first case, or $a,c$ in the second, belong to
one integer-height group in the word with the pair removed.  They can
be made adjacent in either prescribed tie order.  Here is the needed
noninterleaving check even when H is present.  Let $P,Q$ be the
preceding and following occurrences of this rank.  The remaining H
block lies in the cyclic column interval from $P$ to $Q$ which formerly
contained the six-letter-gadget columns.  The two port intervals are
$P\to H_{\rm first}$ and $H_{\rm last}\to Q$.  All other rows at this
integer height are arcs between different consecutive occurrences;
their interiors are disjoint from that interval, or contain it after
choosing a planar cut.  Hence each of the two port rows exchanges with
every such tie row by a strict or shared-endpoint exchange.

Put the input port immediately before the output port in the base
cyclic order.  By the second framing fact move the H internal rows
between them.  This restores \eqref{eq:exact-marked-frame} exactly and
allows the procedure to be repeated for the next inverse pair.

The negative--positive case is obtained by reversing the local row
coordinate in every displayed move, interchanging upper and lower
ranks.  This conjugates the move formulas, not the protected summand:
its internal rows are labelled in the opposite order while describing
the reflected formulas, and those labels are then translated back.
Every actual step is still an ordinary grid move and the protected
marking data are unchanged.  A marked rank not touched by the cancelled letters is an
exterior port: move H away from the local block by the first framing
fact and apply the unmarked cancellation.  Thus every case is covered.
The rank carrying the mark changes by the transposition of the
cancelled letter when the marked connector is involved, exactly as
following that strand through its first crossing prescribes.

In particular, for a factor $\gamma\eta\gamma^{-1}$ with a marked
component, lift the cyclic column rotation moving $\gamma$ from the front to the end
by moving the protected H block aside using its exact frame.  The
word becomes $\eta\gamma^{-1}\gamma$, and the mark is at the interface
of the inverse words.  Cancel the innermost inverse pair, using the
marked rule when its rank meets the mark and the unmarked rule
otherwise.  Reframe after each cancellation.  Induction removes all
of $\gamma^{-1}\gamma$ and leaves H attached to the permutation-tracked
component of $\eta$.  The endpoint is checked by the explicit port
incidences above.  No assertion that unmarked isotopy implies marked
grid equivalence is used.

More explicitly, if $\gamma$ has successive generator indices
$i_1,\ldots,i_m$, the central pairs cancel in that order.  A starting
mark $k$ ends at
$s_{i_m}\cdots s_{i_1}(k)=\pi(\gamma)^{-1}(k)$, with the composition
convention of the brief.  Since
$\pi(\gamma\eta\gamma^{-1})=\pi(\gamma)\pi(\eta)\pi(\gamma)^{-1}$,
membership in the nontrivial cycle of the conjugate implies that this
ending mark belongs to the nontrivial cycle of $\eta$.

In the application the marked component belongs to the nontrivial
permutation cycle of the factor.  Conjugation and the successive
interface cancellations preserve that property for the tracked
component.  Its surviving rank therefore cannot become unused.  Only
an unmarked rank can produce an auxiliary $2\times2$ block in these
steps, exactly as recorded in the two cancellation cases.

All new or removed O markings in a boundary absorption lie on the
marked component.  Shared-endpoint exchanges involve the same
component at their shared corner.  Give a new O the colour of that
component.  Strict exchanges retain every label, and the stabilization
variable identifications only identify equal colours.  Therefore the
entire marked construction is natural for the component-colour
derivative operators by the differentiated domain identities already
proved.

\end{proof}


\section{A signed-factorization scissor construction}

Write $r=N+B$, where $B$ is the total smooth ramification order,
with each smooth local Coxeter word split into its positive-generator factors.  Let $a_l$ be the
length of the nontrivial permutation cycle of the $l$th monodromy factor:
$a_l=a_j$ at a singular point, and $a_l=2$ at a smooth simple branch point.

\begin{lemma}[Exponent count; uses the word lengths and invariant identities]
\label{lem:exponent-count}
\formalname{exteriorRamificationCount}
Suppose the outer event word has an Artin-group factorization
\[
 w_{\rm out}=q_1\cdots q_r,
 \qquad q_l=\gamma_l\eta_l\gamma_l^{-1},
\]
where the $\eta_l$ are the $N$ local positive words and $B$ positive generators.
Then
\[
 B=S-\sum_{j=1}^N a_j+N-1+\mu(V).
\]
\end{lemma}
\begin{proof}
The exponent-sum homomorphism sends an Artin generator to $1$ and is invariant
under conjugation.  Hence the supplied length identities give
\[
 2\delta(W)+S-b
 =\sum_j(2\delta_j+a_j-1)+B.
\]
Substitute $\delta(W)-\sum_j\delta_j=\delta(V)$ and
$2\delta(V)=\mu(V)+b-1$.
\end{proof}

\begin{lemma}[A common strand; new combinatorial]
\label{lem:hurwitz}
\formalname{commonStrandHurwitz}
Let $c_1,\ldots,c_r$ be nontrivial cycles on a finite set, and suppose the group
they generate is transitive.  Fix a letter $o$.  A finite algorithm of Hurwitz
moves transforms the tuple into one in which every cycle moves $o$.
The algorithm lifts to a tuple of arbitrary braid factors having these
permutations, preserving its product and the individual conjugacy classes.
\end{lemma}
\begin{proof}
Call a factor good when its support contains $o$.  At least one factor is good.
If some are bad, transitivity of the support hypergraph provides a bad cycle
$c$ whose support meets that of a good cycle $g$.  Move $c$ next to $g$, keeping
$c$ itself unchanged: moving $c$ past a factor $h$ replaces $h$ by
$chc^{-1}$ or $c^{-1}hc$.  Since $c(o)=o$, this does not change whether $h$ is
good.  If necessary use one further Hurwitz move to put a good factor first.

For $H(g,c)=(gcg^{-1},g)$, direct multiplication gives
\[
 H^{2k}(g,c)=((gc)^k g(gc)^{-k},(gc)^k c(gc)^{-k}).
\]
The orbit of $o$ under $gc$ meets $\operatorname{supp}(g)\cap
\operatorname{supp}(c)$.  Indeed, until that orbit first enters
$\operatorname{supp}(c)$, the map $gc$ acts as $g$.  A full $g$-cycle visits
the nonempty intersection; its first visit to the support of $c$ lies in the
intersection.  Choose an integer $k$ such that $(gc)^{-k}(o)$ belongs to that
intersection.  After $H^{2k}$ both factors are good.  The number of good
factors has increased, so repetition terminates.  All the displayed formulas
are group identities and thus apply to braid factors as well.
\end{proof}

\begin{lemma}[Prescribed factor order; new combinatorial]
\label{lem:ordered-hurwitz}
\formalname{orderedCommonStrandHurwitz}
Suppose every cycle in a tuple of nontrivial cycles moves a distinguished
letter $o$.  Any prescribed order of the labelled conjugacy classes can
be obtained by a finite sequence of Hurwitz moves such that every cycle
in the final tuple still moves $o$.  The same algorithm applies to braid
factors, preserves their product, and keeps each factor in its labelled
conjugacy class.
\end{lemma}
\begin{proof}
It suffices to exchange two adjacent labelled factors $g,c$ while
retaining the distinguished letter in both final supports.  Put $P=gc$.
The previous Hurwitz formula gives
\[
 H^{2k+1}(g,c)=
 \bigl(P^k g c g^{-1}P^{-k},\ P^k gP^{-k}\bigr).
\]
The first entry is a conjugate of $c$ and the second a conjugate of
$g$, so their labels are exchanged.  Both entries move $o$ precisely
when
\[
 x_k=P^{-k}(o)\in \operatorname{supp}(g)\cap
                              g(\operatorname{supp}(c)).
\]
There is a nonnegative $k$ with this property.  Start with
$x_0=o\in\operatorname{supp}(g)$.  As long as
$x_k\notin g(\operatorname{supp}(c))$, the point $g^{-1}(x_k)$ is
fixed by $c$, and hence
\[
 x_{k+1}=c^{-1}g^{-1}(x_k)=g^{-1}(x_k).
\]
Thus, until the desired intersection is reached, the sequence follows
the backward $g$-cycle through $o$.  That cycle reaches $g(o)$ in at
most $|\operatorname{supp}(g)|-1$ steps, and
$g(o)\in\operatorname{supp}(g)\cap g(\operatorname{supp}(c))$
because $c$ moves $o$.  This proves termination and gives an explicit
finite search for $k$.

Use this adjacent exchange in an ordinary bubble sort of the factor
labels.  After each completed exchange all factors again move $o$;
the number of label inversions strictly decreases.  The other factors
are untouched, and the displayed Hurwitz identities preserve the
product at every step.  All identities hold in the braid group as
well, so the same signed-word operations prove the braid assertion.
In the application, order the singular factors by $Q_1,\ldots,Q_N$
first, followed by the simple factors in any fixed order.
\end{proof}

\begin{lemma}[Star scissors and the incidence forest; new combinatorial]
\label{lem:scissor-forest}
\formalname{scissorIncidenceForest}
Suppose $r\ge1$ factors on $S$ strands are conjugates of words with
one nontrivial cycle of length $a_l\ge2$, their permutations are
transitive, and every such cycle contains strand $0$.  Pad each
factor as specified below.  The star rule gives
$s=(S-1)(r-1)$ adjacent-row scissors and a final grid with one main
component and $U=\sum_l(S-a_l)$ auxiliary components.  Exactly $U$
specified scissors can use the O-switch map followed by a strict
row commutation; all others use the X identity.  The endpoints
are the same for either choice.  The incidence matrix of these
$U$ O maps, with the main-component row removed, has determinant
one.  Moreover
$s-U=\sum_l a_l-r-S+1$.
\end{lemma}
\begin{proof}
\subsection{Signed word grids and the scissor switches}

Use the signed three-column convention of Section~\ref{sec:conventions}.  For both signs the
integer-rank incoming corner precedes the outgoing corner of that occurrence.
Consequently, at any fixed integer rank, the arcs between successive
occurrences have pairwise disjoint column intervals, except for the unique
closing arc, whose interval contains all the others.  Thus all tie orders at
an integer rank are related by ordinary, four-distinct-endpoint row
commutations.  Half-integer arcs of distinct letters have disjoint intervals.

Pad each factor on both ends by an identity word $\rho\rho^{-1}$, where
$\rho=\sigma_0\cdots\sigma_{S-2}$.  Every rank then occurs, every factor
begins with $+\sigma_0$, and every factor ends with $\sigma_0^{-1}$.
Padding does not alter the factor or its permutation.

Fix a rank $k\geq1$.  Let $p_l$ be the last outgoing $k$-corner of factor $l$
and $q_l$ its first incoming $k$-corner.  In the concatenation, let $R_l$ be
the row with $O$ at $p_l$ and $X$ at $q_{l+1}$, indices taken modulo $r$.
Perform the star sequence
\[
 (R_1,R_2),\ (R_1,R_3),\ldots,(R_1,R_r),
\]
where a pair means exchanging its two X entries.  After the $l$th stage,
factors $2,\ldots,l$ close separately at this rank, while the row $R_1$
connects $p_1$ to $q_{l+1}$.  At the last stage factor $1$ closes as well.
All row intervals at this rank remain pairwise disjoint or nested: the
active interval spans the consecutive factors already processed, the newly
closed factor intervals lie inside it, the remaining short connectors are
disjoint, and the original last connector encloses the rest.  Therefore each
requested pair can be brought to adjacent rows by ordinary row commutations.
None of the swaps creates an $O=X$ collision, since the incoming and outgoing
corners of an occurrence are distinct columns.

Leave rank $0$ unchanged.  The total number of switches is
\[
 s=(S-1)(r-1).
\]
An adjacent-row X switch has the same explicit chain map and filtration bound
as an adjacent-column X switch: the map is the identity on the O-complex,
and transposing both coordinates in the $J$-formula gives the same
$M_X$ difference $\pm1$.  This proves shifts $(0,1/2)$ and degree $0$.


\paragraph{The O replacement at a specified scissor.}
Bring $R_1,R_l$ into adjacent row positions $u,u+1$, in that order.
Let their O columns be $a=p_1<b=p_l$, and let their X columns before
the switch be the current $q_l,q_{l+1}$.  All four columns are
distinct, because every factor has the stated padding.  Transpose
the adjacent-column O formula.  It is the explicit state map
\[
 F_{\rm row}(x)=
 T^{\,\mathbf1_{\{a<x^{-1}(u+1)\le b\}}}x
\]
to the grid obtained by exchanging the O entries of these rows.
It has degree $-1$ and shifts $(0,-1/2)$.  In that grid the two
row intervals are $p_l\!-\!q_l$ and $p_1\!-\!q_{l+1}$.
They are disjoint or nested, with four distinct endpoints: the
first closes factor $l$ and the second spans the active sequence
of factors.  Commute the two entire rows.  Their O entries and
labels return to their original outgoing factor columns, while
their X entries have been exchanged.  This is exactly the X-scissor
endpoint.  The row-commutation map after $F_{\rm row}$ gives the
selected map.  The detecting $K$ of Lemma~\ref{lem:O-calculus},
followed by that same commutation, retains its canonical-class
and derivative identities.  Thus the incidence formula applies
without any new local map assertion.

\subsection{The explicit forest for O-switch selection}

After common-strand normalization, every factor has just one nontrivial
cycle, and that cycle contains $0$.  The unchanged rank-$0$ connections join
all these components into one main component.  Every fixed rank of a factor
gives one auxiliary component, reduced below to an explicit unknot block.  Hence their number is
\[
 U=\sum_{j=1}^N(S-a_j)+B(S-2).
\]
At rank $k\geq1$, the scissor edges form the star with center factor $1$ and
leaves $2,\ldots,r$.  Its vertex at factor $l$ is main if $k$ belongs to that
factor's nontrivial cycle, and is an auxiliary component otherwise.
Transitivity ensures at least one main vertex at each rank.

If factor $1$ is main, select the edges to precisely the auxiliary leaves.
If it is auxiliary, let $l_*$ be the first main leaf; select $(1,l_*)$ and
the edges from $1$ to every other auxiliary leaf.  This is a tree rooted at
the main component, with exactly one edge per auxiliary vertex.  Taking all
ranks gives $U$ selected edges.  The reduced incidence matrix is square with
determinant $1$ over $\mathbb F_2$: remove leaves successively, expanding
along their unique incident edge.  In particular this selection is suitable
for an incidence-determinant formulation of the O-switch/death class
calculation.

Finally,
\[
 s-U=(S-1)(N+B-1)-\sum_j(S-a_j)-B(S-2)
     =\sum_j a_j+B-N-S+1=\mu(V).
\]



\end{proof}


\section{Exact seam coordinates and successive factor cleanup}

\begin{lemma}[Exact seam coordinates and factor cleanup; new combinatorial]
\label{lem:seam-cleanup}
\formalname{cyclicSeamCoordinates}\formalname{seamFactorCleanup}
Assume the scissor factors are $\gamma_l\eta_l\gamma_l^{-1}$, where
$\eta_l$ is supported on its $a_l$ active ranks, contains every
active generator, and has one full $a_l$-cycle, as in the analytic
factorization.  The preparation at a specified arc below gives a triply prepared
grid by one X stabilization and cyclic permutations.  On triply
prepared factors the seam operation is strictly associative, and
the displayed cyclic block rule exposes any chosen factor by
explicit ordinary moves.  The raw scissor endpoint is that cyclic
block grid.  Successive marked cancellations preserve an exact
iterated-seam invariant and reduce all factors to framed core word
grids with their unused-rank blocks.  Applying the proved block
normalizations and singleton contractions deletes exactly those
auxiliary rows and columns and leaves the exact seam of the
framed active core grids.
\end{lemma}
\begin{proof}

Call a grid $D$ of size $n\geq2$ \emph{triply prepared} if
\begin{equation}\label{eq:triple-preparation}
 X_D(0)=0,\qquad O_D(n-1)=0,\qquad X_D(n-1)=n-1.
\end{equation}
These conditions imply all the preparation conditions required of both
the left and the right factor of the supplied binary seam.  In this
section every iterated seam of triply prepared factors uses the identity
row rotation at every stage.

\subsection{Preparation at a specified arc}

Let an arc of a grid $G$ of size $g\geq2$ run from $O(P,b)$ to
$X(Z,b)$.  X-stabilize at the latter marking.  Insert a column $W$
immediately before $Z$ in the column circle, and insert a row $t$
immediately after $b$ in the row circle, with the three markings
\begin{equation}\label{eq:specified-frame-stabilization}
 O(W,t),\qquad X(W,b),\qquad X(Z,t),
\end{equation}
the old $X(Z,b)$ having been removed.  Cut the column circle immediately
before $Z$, and cut the row circle immediately before $t$.  Then $Z$ is
first, $W$ last, $t$ bottom, and $b$ top.  The resulting grid, denoted
$D(G,b)$, is triply prepared and has size $g+1$.  Formula
\eqref{eq:specified-frame-stabilization} is an ordinary X stabilization;
the two choices of cuts are cyclic permutations.  The new O has the
colour of the component containing the specified arc.  No component is
created or merged.

\subsection{Strict associativity and a cyclic block formula}

Let $D_1,\ldots,D_r$ be triply prepared, with sizes $n_l$, and set
\[
 e_l=n_l-1,\qquad c_l=\sum_{i<l}e_i,\qquad d=\sum_{l=1}^r e_l.
\]
Their iterated seam $T$ has size $d+1$.  Substitution in the binary
seam formula shows that $T$ is itself triply prepared.  It also shows
that the seam operation is strictly associative on these data: both
parenthesizations have the same O and X in every column, with offsets
given by the corresponding sums of the $e_l$.  In particular, no
unmentioned grid move or row rotation is concealed in reassociation.
Explicitly, for $0\leq j<e_l$ and $Y\in\{O,X\}$,
\[
 Y_T(c_l+j)=c_l+Y_{D_l}(j),\qquad O_T(d)=0,\quad X_T(d)=d.
\]
This formula follows by induction from the binary formula, because
the unique zero-row O of each right factor lies in its last column
and its last X is in its top row.  It proves the associativity and
preparation assertions directly.

Suppose $r\geq2$.  Let $M=d$ be the last column of $T$ and $A=0$ its
first column.  Across the toroidal last/first column boundary and the
top/bottom row boundary there are exactly the three markings
\[
 X(M,d),\qquad O(M,0),\qquad X(A,0).
\]
The fourth corner $(A,d)$ is empty.  Indeed, the X of $A$ is in row
$0$, while its O is in the first factor's row block, hence in a row
at most $e_1<d$.  X-destabilize this block, deleting $M$ and row $0$
and replacing its X markings by $X(A,d)$ before renumbering.  A cyclic
cut makes this an ordinary adjacent-row, adjacent-column
destabilization if the convention requires a planar $2\times2$ block.

Denote the resulting size-$d$ grid by $C(D_1,\ldots,D_r)$.  Its columns
and rows are divided into blocks of sizes $e_l$.  For
$0\leq j<e_l$ its complete marking formula is
\begin{equation}\label{eq:cyclic-block-grid}
\begin{split}
 O_C(c_l+j)&=c_l+O_{D_l}(j)-1,\\
 X_C(c_l+j)&=
 \begin{cases}
 c_l-1\pmod d,&j=0,\\
 c_l+X_{D_l}(j)-1,&j>0.
 \end{cases}
\end{split}
\end{equation}
All values on the first line lie in the $l$th row block, since the only
O in row $0$ of $D_l$ is in its omitted last column.  The same is true
of the second line for $j>0$.  The first X of a block lies in the top
row of the preceding block, with the blocks read cyclically.  This is
the entire effect of the destabilization.

The formula is literally cyclically symmetric.  To put factor $D_l$
first, cyclically rotate the column blocks by $c_l$ and the row blocks
by the same amount.  These are permitted cyclic column and row
permutations.  Apply the inverse X stabilization at the first-column
X in the top row: add a last column $M$ and a bottom row $0$, place
$O(M,0),X(M,\mathrm{top}),X(0,0)$, and remove the old
first-column/top-row X.  One obtains exactly
\[
 \Sigma(D_l,D_{l+1},\ldots,D_r,D_1,\ldots,D_{l-1}),
\]
where the displayed expression is strictly associative.  Thus any
factor can be exposed as the left factor of a binary seam whose right
factor is the contiguous seam of all the others.  This conclusion is
proved by marking coordinates, without a commutativity assertion for
connected sums.

The formula includes factors of size $2$: then $e_l=1$ and the
$j>0$ case for that block is simply absent.  Since $r\geq2$, $d\geq2$,
so the temporary destabilized grid never has size $1$.  When $r=1$
there is no cyclic factor permutation to perform, and the operation is
omitted; its fourth corner need not be empty in that case.

\subsection{The raw scissor endpoint is already this cyclic grid}

Let $q_1,\ldots,q_r$ be the padded signed factors in the scissor
construction, and let $G_l$ be their standalone word grids.  At each
factor the first letter is $\sigma_0$ and the last letter is
$\sigma_0^{-1}$.  Thus its first column $A_l$ is the first incoming
rank-$0$ corner, its last column $B_l$ is the last outgoing rank-$0$
corner, and its closing rank-$0$ row $b_l$ runs
$O(B_l,b_l)\to X(A_l,b_l)$.

Apply the preparation formula
\eqref{eq:specified-frame-stabilization} to the standalone grid $G_l$
at this arc and call the result $D_l$.  This defines $D_l$ explicitly;
it is not an assertion that unrelated factor moves lift through a
seam.  Its nonlast columns are exactly the columns of $G_l$, in their
original cyclic order, and its nonbottom rows are exactly the rows of
$G_l$, cut just after $b_l$, so that $b_l$ is last.  Its omitted last
column and bottom row are precisely the added frame.

After the scissor switches at every rank other than $0$, the raw
endpoint has these same column blocks.  Every row is internal to one
block except the rank-$0$ rows
\[
 R_l:O(B_l)\longrightarrow X(A_{l+1}),
 \qquad l\pmod r.
\]
Assign $R_l$ to row block $l$.  Within each block put the rows in the
cyclic order of $G_l$ starting immediately after its closing rank-$0$
row; put $R_l$ last.  These row permutations are realized as follows.
Rows internal to different column blocks have disjoint intervals.
For $l<r$, $R_l$ has consecutive column endpoints $B_l,A_{l+1}$;
$R_r$ is the closing interval containing every other column.  Hence
every pair of rows with different block labels has disjoint or nested
intervals, or shares an endpoint.  Every needed adjacent exchange is
therefore either a strict noninterleaving commutation or the explicit
shared-endpoint replacement already proved.

To implement independent cyclic cuts, gather one row colour into a
consecutive block by these exchanges.  Move its first row globally
around the row circle, then exchange it back past all other colours
until it is last in its own block.  Repeating realizes any cyclic
rotation of that colour while preserving all other cyclic orders.
After choosing the required cyclic order for every colour, use a
stable sort of the colour labels, which exchanges only different
colours.  This is a finite algorithm producing the stated row order.

The raw scissor endpoint in this order is exactly
$C(D_1,\ldots,D_r)$ by \eqref{eq:cyclic-block-grid}: its only changed X
is the first one in each block, now on the preceding block's last
row.  Consequently a single global X stabilization, the inverse of
the destabilization above, gives the exact iterated seam of the
triply prepared $D_l$.  This preparation changes neither the saddle
count nor any component colour.  The added O is on the main component.
For $r=1$ the raw endpoint is simply $G_1$, and this same stated
stabilization is exactly $G_1\to D(G_1,b_1)$.  It does not rely on a
preceding cyclic-seam destabilization.

\subsection{The invariant under marked inverse-pair cancellation}

Maintain the following exact invariant: the current grid is the
strict iterated seam of triply prepared factor grids, one for each
monodromy factor; a factor consists of its current marked word part
and the explicitly retained $2\times2$ blocks of its unused ranks,
each with its two rows and its two columns adjacent.
All such auxiliary components keep their prescribed colours.  The
cyclic block rule exposes any desired factor $D$ first.  By strict
associativity the resulting grid is exactly
\[
 \Sigma(D,H),
\]
where $H$ is the contiguous seam of the other factors.  For $r\geq2$,
$H$ has size at least $2$, is triply prepared, and its complete
marking data will be protected during the following operation.

Here $D=D(G,b)$ is the one-stabilization preparation of its underlying
marked word part, including any retained auxiliary blocks.  The binary
seam formula already gives its exact framed insertion form: the word
column $Z$ is retained with X in the bottom row $t$, the first H
column has X in the top base row $b$, and the last H column has O in
$t$.  Its internal H rows lie above $b$.  No additional stabilization
to restore the virtual frame column is performed.  Cancel the padding and the
conjugator of the exposed factor by the marked inverse-pair theorem, using unmarked
inverse-pair cancellation whenever the current mark does not meet the
two affected ranks.  At every marked cancellation the protected H
block retains its columns, internal rows, their orders, and every
internal marking.  Only its incoming and outgoing port rows change.
The local theorem restores the exact frame after each cancellation.
For a factor $\gamma\eta\gamma^{-1}$, the successive transpositions
track the mark into the nontrivial permutation cycle of $\eta$.
Therefore the marked rank never becomes unused.  The unused ranks are
exactly the $S-a_l$ fixed strands of that factor, and each is retained
as one explicit isolated $2\times2$ grid.

It remains to check that the output frame recovers the same global
invariant, rather than merely representing an attached summand.  Let
the resulting base word part have the arc
$O(P,b)\to X(Z,b)$.  Its framed insertion has H input
$X(H_{\rm first},b)$ and H output $O(H_{\rm last},t)$, while the
retained base column $Z$ has X in row $t$.  Introduce the
\emph{virtual} last column $W$ with $O(W,t),X(W,b)$.  First use the
frame's proved column-exchange rule to move the entire H block
immediately before the retained base column $Z$, preserving all
internal and exterior column orders.  Then cut immediately before
$Z$, so the base columns come first and the H block last.  With row
cut $t$ first and $b$ last among the base rows, the
virtual base grid is exactly $D(G,b)$ from
\eqref{eq:specified-frame-stabilization}.  The virtual column is
precisely the column deleted by the seam formula.  Substitution in
that formula recovers all actual markings of the frame, including
$X(H_{\rm first},b)$ and $O(H_{\rm last},t)$.  Thus the actual output
grid is literally
\[
 \Sigma(D(G,b),H).
\]
The exposed factor has been replaced by another triply prepared one,
and every other factor is unchanged.  This proves the induction
invariant and permits the next factor to be exposed by the cyclic
block rule.  No general statement that an unmarked equivalence lifts
through a seam is needed.

Postpone all auxiliary contractions until every factor has been
cleaned.  An auxiliary $2\times2$ block can then be placed at the
prescribed location by ordinary moves.  First move its two columns together
to the desired column position: their row endpoints are the same two
adjacent rows, and no exterior column uses either row, so all these
exchanges are strictly noninterleaving.  With its columns adjacent,
move its two rows to the desired row position; the same argument
applies with coordinates transposed.  The block has no shared marking
with the main component.  If necessary exchange its two rows by the
proved shared-endpoint rule to give the four-mark pattern of
Proposition~\ref{prop:four-mark-collapse}.

For each block first apply that proposition's explicit degree-zero,
shift-zero collapse to a doubled singleton.  These maps preserve the
canonical class and intertwine its component-colour derivative.  Only
after these collapses apply the degree-$+1$, shift-$(0,1/2)$ singleton
contractions of Lemma~\ref{lem:dummy-transport}, using the persistent
main component as root.  These are the $U$ contractions in the
incidence-determinant calculation.  A four-mark block is never
treated as an ordinary three-corner destabilization.  Intermediate
doubled singletons are auxiliary marking pairs with the explicitly
defined complexes and maps of those two results.

Their exact target formulas delete only the selected auxiliary rows
and columns, retaining all others in their original relative order.
In particular the first-column X, bottom O, and last-column top X of
every main frame remain.  At the end the actual grid is exactly the
seam of the same factor frames with their auxiliary blocks removed.
The persistent main component has a row with distinct O and X
throughout, as required for the contractions.  This proves the target
claim at the level of marking permutations as well as the required
colour naturality.

After all factors have been processed and the prescribed $U$ blocks
have died, each remaining factor is the frame preparation of
$\operatorname{wordGrid}(\eta_l)$ at a rank in its unique active
cycle.  Its size is at least $3$, since the unframed word grid has size
at least $2$.  Reversing its one frame stabilization, together with
its explicit cyclic cuts, gives an ordinary-move equivalence with
$\operatorname{wordGrid}(\eta_l)$.  Thus the preparation condition and
the required \texttt{MovesTo} condition are both verified without any
further rank normalization.


\end{proof}


\section{Removing the simple factors and fixing the target}

\begin{lemma}[Simple-factor projection and exact target;
new combinatorial, uses the binary seam comparison]
\label{lem:simple-factor}
\formalname{simpleFactorProjection}
Suppose the cleaned factors have been ordered with the singular
factors first in their prescribed order and the simple factors last.
The functional \eqref{eq:simple-factor-functional} on
$C(\operatorname{wordGrid}(\sigma_0))$ is a $\Lambda$-linear
degree-zero chain map to $\Lambda$ with shifts $(0,0)$ and sends
the canonical class to $1$.  Consequently the seam of the cleaned
singular and simple factors admits a degree-zero, shift-zero,
canonical-class-preserving map to the prescribed iterated seam
of the singular factors alone.
\end{lemma}
\begin{proof}

All ordinary moves and all prescribed auxiliary deaths in the preceding
section occur before using a seam--tensor comparison.  Thus the
component-colour derivative proof of the O-switch/death canonical-class
identity applies to the explicit marking pairs and their specified maps
throughout.  At that stage the canonical-class
calculation is the one completed in Section~\ref{sec:final-proof}.
The following operations add no further degree or filtration shift.

For a smooth-ramification factor, its active word is a positive
generator on two strands.  Its three-by-three word grid has
\[
 O=(1,0,2),\qquad X=(0,2,1).
\]
Here a state such as $021$ denotes the values of its row permutation
in increasing column order.  Direct enumeration of the two possible
empty rectangles for each transposition gives
\[
\begin{array}{c|c}
 x&\partial x\\\hline
 012&021+102+T\,210\\
 120&021+102+T\,210\\
 201&021+102+T\,210\\
 021&0\\
 102&0\\
 210&0.
\end{array}
\]
The Maslov and Alexander values of the states used below are
\[
 (M_O,A)(021)=(-2,-1),\qquad
 (M_O,A)(102)=(-2,-2),\qquad
 (M_O,A)(210)=(0,0).
\]
With $\Lambda$ having zero differential and its usual monomial degree
and bifiltration, define the $\Lambda$-linear map
\begin{equation}\label{eq:simple-factor-functional}
 \epsilon(210)=1,\qquad \epsilon(021)=T,\qquad
 \epsilon(012)=\epsilon(102)=\epsilon(120)=\epsilon(201)=0.
\end{equation}
The displayed differential table gives $\epsilon\partial=T+T=0$.
The two nonzero images have exactly the degrees of their inputs.
Their filtration coordinates satisfy
\[
 (0,0)=(0,A(210)),\qquad (-1,-1)\leq(0,A(021))=(0,-1).
\]
Multiplication by $T^k$ translates both input and output coordinates
by $(-k,-k)$, proving shifts $(0,0)$ for every monomial.  The degree-zero
cycle $210$ lies in the polynomial subcomplex, and
\eqref{eq:simple-factor-functional} proves its class nonzero.  Hence
it is the supplied canonical class and is sent to $1$.

Before the scissor construction, use
Lemma~\ref{lem:ordered-hurwitz} to order the labelled singular factors
as $Q_1,\ldots,Q_N$, followed by the $B$ simple factors, while keeping
strand $0$ in every nontrivial factor cycle.  The subsequent factor
exposures change this order only by cyclic rotation: each uses the
cyclic block rule, and marked cancellation changes only the exposed
factor.  After cleanup and the auxiliary contractions, apply that same
cyclic block rule to put the factor labelled $Q_1$ first.  This is
possible because $N\ge1$.  The actual grid is now the exact seam of
the cleaned singular factors in their prescribed order, followed by
the cleaned simple factors.  All these ordering operations are
ordinary moves of degree and shift zero.

For a cleaned simple factor $H$, define
\[
 \epsilon_H=\epsilon\circ M_H:C(H)\longrightarrow\Lambda,
\]
where $M_H$ reverses its one frame stabilization and cyclic cuts to
the three-by-three word grid.  This map is $\Lambda$-linear.  Indeed,
the ordinary maps chosen here are the explicitly proved polynomial
domain maps: a commutation is its pentagon count, the primitive
stabilization directions are the incoming and outgoing domain counts,
and the other corners use their finite corner reductions.  Each is a
sum of state terms with powers of $T$, extended $\Lambda$-linearly.
Cyclic row and column permutations use their literal state bijections
and preserve rectangle weights; these too are $\Lambda$-linear.
The earlier proofs give their degree, shift, and class properties.
Consequently $\epsilon_H$ has degree zero, shifts $(0,0)$, and sends
$\operatorname{can}(H)$ to $1$.

Remove the last simple factor one at a time.  By the strict seam
formula the current grid is $\Sigma(G,H)$, where $H$ is that last
simple factor and $G$ is the exact seam of the remaining prefix.
The prefix always includes the $N\ge1$ singular factors, so it is a
valid grid of size at least two.  Apply the supplied \emph{binary}
comparison, followed by
\begin{equation}\label{eq:ordered-binary-stripping}
 C(\Sigma(G,H))\xrightarrow{\ \Theta_{G,H}\ }
 C(G)\otimes_\Lambda C(H)
 \xrightarrow{\ \operatorname{id}\otimes\epsilon_H\ }
 C(G)\otimes_\Lambda\Lambda\cong C(G).
\end{equation}
The second arrow is well defined because $\epsilon_H$ is
$\Lambda$-linear.  It is a chain map of degree zero and shifts
$(0,0)$: if a monomial term of $\epsilon_H(y)$ is $T^m$, its shift
bounds say $m\ge0$ and $-m\le A_H(y)$, so the image of
$T^k(x\otimes y)$ has coordinates
\[
 (-k-m,A_G(x)-k-m)
 \le (-k,A_G(x)+A_H(y)-k).
\]
The canonical-class identity of the binary comparison, followed by
$\epsilon_H\operatorname{can}(H)=1$, proves the identity for
\eqref{eq:ordered-binary-stripping}.

Repeat this operation exactly $B$ times.  Its exact target after each
step is the prefix grid $G$, so no comparison is tensored with an
identity map and no inverse or iterated tensor comparison is used.
In particular no $\Lambda$-linearity of the supplied
$\Theta_{G,H}$ is assumed.  The final target is literally the seam
of the retained singular factors $D_1,\ldots,D_N$ in their required
order, with identity row rotations.  Every step has degree zero,
shifts $(0,0)$, and preserves the canonical class.  If $B=0$ this
entire removal operation is the identity.

The original scissor uses $s$ saddles, of which $U$ are O switches;
the other $s-U$ are X switches.  The O switches and their $U$ deaths
cancel in degree and filtration shift.  Since the exact count proved
earlier gives $s-U=\mu(V)$, the resulting map has degree zero and
shifts $(0,\mu(V)/2)$.  All size bounds required in Statement E hold,
and every local factor is related by ordinary moves to its supplied
positive word grid.

If $r=1$, transitivity forces the sole nontrivial cycle to contain all
$S$ ranks.  Since there is at least one singularity, this sole factor
is a singularity factor.  Thus $U=s=0$ and the exact count gives
$\mu(V)=0$.  There is no scissor and no factor permutation.  Apply
standalone cyclic word rotation and the unmarked free-cancellation
rule to reach its local positive word grid, then the one stabilization
\eqref{eq:specified-frame-stabilization} gives its prepared target.
This separately covers the only case in which the cyclic seam
destabilization was not used.

\end{proof}


\section{Construction of the complete map and proof of Statement E}
\label{sec:final-proof}

\begin{proof}[Proof of Theorem~\ref{thm:E}]
Choose the assembly supplied by the brief.  Apply the finite-projection
and event-factorization lemma to choose a sufficiently small allowed
level and a projection transverse at every central and local branch.
It gives the central positive word
$\alpha=\mathbf d_S\mathbf w$, the local positive words $\beta_j$,
and a finite derivation, using signed Artin relations and conjugacy,
of
\[
 \alpha\ \sim\ q_1\cdots q_r,\qquad
 q_l=\gamma_l\eta_l\gamma_l^{-1},\qquad r=N+B.
\]
If the initial comparison is by whole-word conjugacy, absorb its
conjugator $c$ into every $\gamma_l$.  Indeed
$c(q_1\cdots q_r)c^{-1}=\prod_l(cq_lc^{-1})$, with a finite free
derivation obtained by inserting $c^{-1}c$ between consecutive
factors.  The guarded word normalization can therefore use only
local Artin and free relations in its suffix.
Each $\eta_l$ is a local $\beta_j$, supported on its first $a_j$
strands, or $\sigma_0$ on its first two strands.  A cluster originally
in another consecutive interval is put in this form by the explicit
identity $\operatorname{sh}_k(\eta)=\rho_S^k\eta\rho_S^{-k}$.
The factor permutations generate a transitive group on $S$ points.
In particular every local degree is at most $S$, since it is a
cluster in an $S$-point regular fibre.

Apply the finite Hurwitz algorithm with distinguished strand $0$.
It preserves the product and the individual conjugacy classes and
puts strand $0$ in every nontrivial factor cycle.  Remember the
original singularity index attached to each local factor.  The
odd-Hurwitz interchange rule then orders the factors as the singular
factors $1,\ldots,N$, followed by all the simple factors, with strand
$0$ still in every nontrivial cycle after each completed interchange.
Pad each
factor on both ends by
$p=\rho_S\rho_S^{-1}$.  Since $p^{-1}=p$ as a literal signed word,
\[
 p\gamma_l\eta_l\gamma_l^{-1}p
       =(p\gamma_l)\eta_l(p\gamma_l)^{-1}.
\]
The signed-word rules, with the external guard described there,
give an explicit ordinary-move map from the specified native source
grid to the grid of this concatenation.  It has degree and shift zero
and sends the canonical class to the canonical class.

Perform the star scissors at ranks $1,\ldots,S-1$.  There are
$s=(S-1)(r-1)$ operations.  Select the $U$ edges of the explicitly
specified rooted forest, where
$U=\sum_j(S-a_j)+B(S-2)$.  At every unselected edge use the
adjacent-row X identity.  At every selected edge use the transposed
diagonal O-switch map followed by the specified noninterleaving
whole-row commutation.  Thus both choices have exactly the same
endpoint markings.  The selected choice has degree $-1$ and shifts
$(0,-1/2)$, while the unselected choice has degree zero and shifts
$(0,1/2)$.

Colour the O markings by their components at the completed scissor
endpoint: colour $0$ is the single main component, and the other
$U$ colours are its auxiliary components.  Transport these labels
back through the scissors.  This is legitimate without any
monochromatic assumption on an earlier component.  Every intervening
row exchange during the scissors has four distinct endpoint columns,
so it uses a labelled commutation and identifies no O variables.
An X switch leaves the O differential and its labels unchanged.
For a selected O switch the two final endpoint colours are exactly
the two vertices of its forest edge.  The final incidence matrix,
with its root row removed, therefore has determinant one.

Use the explicit cyclic-grid stabilization to obtain the exact seam
of triply prepared factors.  Expose one factor at a time by the
cyclic seam formula.  Its padding and conjugator are removed by the
marked inverse-pair rules while the complementary factor grid is
retained as the protected block.  The exact frame is restored after
every cancellation and the resulting factor is again triply
prepared.  Thus this is an actual finite sequence on the whole grid.
The tracked rank is in the core's nontrivial cycle.  Precisely its
$S-a_l$ fixed ranks become four-mark $2\times2$ components; the marked
rank never disappears.  Complete this procedure for every factor.

For a retained four-mark block with the opposite O diagonal, exchange
its two adjacent rows by the explicit shared-endpoint rule, so that
it has the displayed collapse pattern.  Normalize each of the $U$
four-mark blocks to a doubled singleton by
the explicit degree-zero, shift-zero, colour-natural map of
Proposition~\ref{prop:four-mark-collapse}.  These are algebraic
normalizations, with their full formulas proved above.  They cannot
be replaced by three-corner destabilizations.  All ordinary moves
and these normalizations preserve the canonical class and intertwine
every component-colour derivative on homology.  Denote their
composite, following the scissors, by $P$.

Finally remove the $U$ doubled singletons by the explicit contraction
maps, always taking a surviving O of the main component as root.
For the singleton of colour $c$ its contraction is, on homology,
$\rho_c\Psi_c$, with $\rho_c$ of degree zero and second shift $1/2$,
preserving the canonical class.  A derivative of another dying
colour commutes past $\rho_c$; the possible root-substitution
correction contains $\Psi_c^2$ and is zero.  Hence the induced map of
all contractions is the ordered product of the $\rho_c$'s following
the product of all $U$ dying-colour derivatives.  Pass those
derivatives back through $P$ and use the determinant identity
\eqref{eq:determinant-pairing}.  Since the reduced forest incidence
matrix has determinant one, the result before the $\rho_c$'s is
exactly the canonical class.  Each $\rho_c$ preserves that class.
This proves nonvanishing and, by Lemma~\ref{lem:N}, the required
canonical-class identity for the entire map so far.

The target markings at this stage are literally the seam of the
cleaned triply prepared factor grids.  Deleting a singleton deletes
only its row and column and preserves the order of all retained
rows and columns.  The auxiliary blocks occupy no main frame port,
so these exact deletions commute with the displayed seam marking
formula.  Every cleaned factor is the one-frame preparation of the
word grid of its $\eta_l$, at its tracked active arc.

The cyclic exposures preserve the cyclic order of all factors.
Use the explicit cyclic seam rule to restore the chosen first factor,
so that the singular factors are again $1,\ldots,N$ and the simple
factors are last.  If a simple factor remains, write the actual grid
as $\Sigma(G,H)$, where $H$ is its last prepared simple factor.
Apply the supplied \emph{binary} comparison
\[
 C(\Sigma(G,H))\longrightarrow C(G)\otimes_\Lambda C(H).
\]
Reverse the frame preparation of $H$ and apply
\eqref{eq:simple-factor-functional}, obtaining a $\Lambda$-linear
map $\epsilon_H:C(H)\to\Lambda$ of degree and shift zero.
Consequently $\operatorname{id}_{C(G)}\otimes_\Lambda\epsilon_H$
is well defined and gives a map to $C(G)$.  This composition uses
no $\Lambda$-linearity of the supplied binary comparison.  It
preserves the canonical class.  Repeat for the $B$ final simple
factors.  No iterated tensoring of that comparison, or derivative
naturality for it, is used.

All these maps have degree zero and shifts $(0,0)$.  Put $D_j$ equal
to the retained prepared grid of
$\beta_j$.  It is related to $\operatorname{wordGrid}(\beta_j)$
by the explicitly given moves.  The exact final grid is
$\operatorname{iteratedSeamGrid}(D)$, with every row rotation equal
to the identity.

The degree of the composite is $-U+U=0$.  Its first shift is zero.
Its second shift is
\[
 \frac{s-U}{2}-\frac U2+\frac U2
     =\frac{s-U}{2}
     =\frac{\mu(V)}2.
\]
Every equality uses the explicitly verified bounds for its maps;
in particular the four-mark normalizations contribute zero.
The available local words give $a'_j=a_j\in[2,t_J]$ and all their
required cycle, generator, and length properties.  The central
word has the required properties by (I1)--(I2).

If $r=1$, the separate construction in the endpoint lemma uses
standalone conjugator cancellation and one frame stabilization.
Transitivity gives $a_1=S$, and $N\ge1$ gives $N=1,B=0$.
Thus $U=s=\mu(V)=0$, and the same conclusion follows with no
scissors, auxiliary contractions, or tensor-factor deletions.
This proves every assertion of Statement E.
\end{proof}

\subsection{Complete operation count}

The finite generating rules above specify every marking permutation,
every adjacent pair, and every three-corner destabilization.  Since
this is a direct algebraic construction, its count must distinguish
the maps of nonzero degree from ordinary X switches:
\[
\begin{array}{l|c|c|c}
\text{operation}&\text{number}&\text{degree}
 &\text{second shift}\\\hline
\text{scissor X identity}&s-U=\mu(V)&0&1/2\\
\text{selected scissor O map}&U&-1&-1/2\\
\text{four-mark normalization}&U&0&0\\
\text{final singleton contraction}&U&+1&1/2\\
\text{ordinary and endpoint maps}&\text{finite}&0&0
\end{array}
\]
There are no births.  If every four-mark normalization is expanded
into its defining O switch and singleton contraction, the complete
formula uses $2U$ O switches and $2U$ singleton contractions, besides
the $\mu(V)$ X identities.  The added pairs still have total degree
and shift zero.  Thus the effective X count minus birth count is
exactly $\mu(V)$, and the actual shift computation is the displayed
sum, not an application of the restricted M/X/B movie assertion.

All counts are explicit functions of the brief's invariants:
\[
\begin{split}
B&=S-\sum_j a_j+N-1+2\delta(V)-b+1,\\
r&=N+B,\qquad s=(S-1)(N+B-1),\\
U&=NS-\sum_j a_j+B(S-2).
\end{split}
\]
The number of zero-shift moves is finite by the stated algorithms:
the finite event derivation, finite Hurwitz iteration, finite
bubble sorts, the 29-step Artin rule, and finitely many inverse
cancellations.  No bound on that harmless count is needed.


\section{A fully indexed example with \texorpdfstring{$S=4$ and $\mu=1$}{S=4 and mu=1}}
\label{sec:worked-four}

This is a numerical and combinatorial instance of the construction;
no assertion is made that these data are realized by a holomorphic
counterexample.  Write $012$ for $\sigma_0\sigma_1\sigma_2$, and take
\[
 \alpha=(012)^5 0=\mathbf d_4\,(0120),\qquad
 \beta_1=(012)^5,\qquad \eta_2=\sigma_0.
\]
The numerical data are
\[
 S=4,\quad M=2,\quad b=2,\quad N=1,\quad a_1=t_J=4,
 \quad \delta(V)=1,\quad \delta(W)=7,\quad\delta_1=6,
 \quad\mu(V)=1.
\]
They satisfy $2\delta(V)+M=S$, $\delta(W)=\delta(V)+\delta_1$,
$\mu(V)=2\delta(V)-b+1$, and $\delta(W)\ge S(S-1)/2$.
The lengths are $\ell(\alpha)=16=2\delta(W)+S-b$ and
$\ell(\beta_1)=15=2\delta_1+a_1-1$.  With the stated permutation
convention,
\[
 \pi(\beta_1)=(0\ 1\ 2\ 3),\quad
 \pi(\eta_2)=(0\ 1),\quad
 \pi(\alpha)=(0\ 2\ 3)(1).
\]
Thus $B=1$, both factor cycles contain $0$, and the common-strand
algorithm needs no Hurwitz move in this example.

For the source and the unprepared local grid, order equal-height arcs
by their outgoing column.  The columns in the following tables are
zero-indexed; each pair of rows $O(c),X(c)$ specifies the complete
native tuple $(O,X)$.

\paragraph{Source grid: size $48$.}\mbox{}
\[
\begin{array}{c|rrrrrrrrrrrr}
c&0&1&2&3&4&5&6&7&8&9&10&11\\
O(c)&6&0&12&23&13&28&38&29&43&7&1&14\\
X(c)&5&22&6&12&37&23&28&47&38&0&13&7\\
\hline
c&12&13&14&15&16&17&18&19&20&21&22&23\\
O(c)&24&15&30&39&31&44&8&2&16&25&17&32\\
X(c)&14&29&24&30&43&39&1&15&8&16&31&25\\
\hline
c&24&25&26&27&28&29&30&31&32&33&34&35\\
O(c)&40&33&45&9&3&18&26&19&34&41&35&46\\
X(c)&32&44&40&2&17&9&18&33&26&34&45&41\\
\hline
c&36&37&38&39&40&41&42&43&44&45&46&47\\
O(c)&10&4&20&27&21&36&42&37&47&11&5&22\\
X(c)&3&19&10&20&35&27&36&46&42&4&21&11\\
\end{array}
\]

Put
\[
 p=012\,\overline2\,\overline1\,\overline0,
 \qquad q_1=p\beta_1p,\qquad q_2=p0p,\qquad Q=q_1q_2,
\]
where a bar denotes a negative generator.  Their lengths are
$27,13,40$, respectively.  Inserting the four displayed identity
words gives an explicit free-insertion derivation from $\alpha$ to
$Q$; the signed-word rules realize it by ordinary grid moves.
The first factor uses columns $0,\ldots,80$, and the second uses
columns $81,\ldots,119$.

At each integer rank, put all internal arcs first in outgoing-column
order, then put the two connector rows $R_1,R_2$ last, in that order.
At each half-integer height use outgoing-column order.  These
tie changes are the already proved ordinary row moves.  The resulting
complete padded grid is as follows.

\paragraph{Padded grid $G^0$: size $120$.}\mbox{}
\[
\begin{array}{c|rrrrrrrrrrrr}
c&0&1&2&3&4&5&6&7&8&9&10&11\\
O(c)&14&0&28&55&29&68&94&69&107&95&108&70\\
X(c)&13&54&14&28&93&55&68&119&94&107&69&95\\
\hline
c&12&13&14&15&16&17&18&19&20&21&22&23\\
O(c)&56&71&30&15&31&1&16&2&32&57&33&72\\
X(c)&70&29&56&30&0&15&1&31&16&32&71&57\\
\hline
c&24&25&26&27&28&29&30&31&32&33&34&35\\
O(c)&96&73&109&17&3&34&58&35&74&97&75&110\\
X(c)&72&108&96&2&33&17&34&73&58&74&109&97\\
\hline
c&36&37&38&39&40&41&42&43&44&45&46&47\\
O(c)&18&4&36&59&37&76&98&77&111&19&5&38\\
X(c)&3&35&18&36&75&59&76&110&98&4&37&19\\
\hline
c&48&49&50&51&52&53&54&55&56&57&58&59\\
O(c)&60&39&78&99&79&112&20&6&40&61&41&80\\
X(c)&38&77&60&78&111&99&5&39&20&40&79&61\\
\hline
c&60&61&62&63&64&65&66&67&68&69&70&71\\
O(c)&100&81&113&21&7&42&62&43&82&101&83&114\\
X(c)&80&112&100&6&41&21&42&81&62&82&113&101\\
\hline
c&72&73&74&75&76&77&78&79&80&81&82&83\\
O(c)&102&118&84&63&92&44&22&53&12&23&8&45\\
X(c)&114&83&102&84&43&63&44&7&22&12&53&23\\
\hline
c&84&85&86&87&88&89&90&91&92&93&94&95\\
O(c)&64&46&85&103&86&115&104&116&87&65&88&47\\
X(c)&45&92&64&85&118&103&115&86&104&87&46&65\\
\hline
c&96&97&98&99&100&101&102&103&104&105&106&107\\
O(c)&24&48&9&25&10&49&26&11&50&66&51&89\\
X(c)&47&8&24&9&48&25&10&49&26&50&88&66\\
\hline
c&108&109&110&111&112&113&114&115&116&117&118&119\\
O(c)&105&90&117&106&119&91&67&93&52&27&54&13\\
X(c)&89&116&105&117&90&106&91&51&67&52&11&27\\
\end{array}
\]

The three scissor operations are now adjacent-row operations without
any preliminary row exchanges.  Both their row positions and the
columns containing the exchanged entries are explicit:
\[
\begin{array}{c|c|c|c|c}
\text{rank}&(R_1,R_2)&\text{O columns}&\text{X columns}&\text{operation}\\\hline
1&(53,54)&(79,118)&(1,82)&X\\
2&(92,93)&(76,115)&(4,85)&O\text{, then whole-row commutation}\\
3&(118,119)&(73,112)&(7,88)&O\text{, then whole-row commutation}
\end{array}
\]
Here the X operation is the explicitly proved adjacent-row version:
its chain map is the identity on the O complex.  The listed X columns
identify the two tuple entries; they are not asserted to be adjacent
columns.  Transposing the adjacent-column grading calculation proves
its second shift $1/2$.

For exact intermediate tuples, let $\tau_{a,b}Y$ mean exchange entries
$a,b$ of a permutation tuple $Y$, and let $s_{u,v}Y$ mean exchange
the row values $u,v$ wherever they occur.  If the displayed padded tuple
is $(O,X)$, put
\[
 X_1=\tau_{1,82}X,\qquad
 X_2=\tau_{4,85}X_1,\qquad
 X_3=\tau_{7,88}X_2.
\]
The full successive grid list, including the O-switch targets, is
\[
\begin{split}
G^0&=(O,X),\\
G^1&=(O,X_1),\\
H^2&=(\tau_{76,115}O,X_1),\\
G^2&=(s_{92,93}\tau_{76,115}O,s_{92,93}X_1)=(O,X_2),\\
H^3&=(\tau_{73,112}O,X_2),\\
G^3&=(s_{118,119}\tau_{73,112}O,s_{118,119}X_2)=(O,X_3).
\end{split}
\]
Every tuple is valid.  The two whole-row commutations are strict:
after the first O exchange their marking intervals are $[4,76]$ and
$[85,115]$; after the second they are $[7,73]$ and $[88,112]$.
Thus no shared-endpoint convention is needed here.  On a state $x$
the O-switch maps in these coordinates are explicitly
\[
 F_2(x)=T^{\mathbf1_{\{76<x^{-1}(93)\le115\}}}x,
 \qquad
 F_3(x)=T^{\mathbf1_{\{73<x^{-1}(119)\le112\}}}x.
\]
These are the transposed O-switch formulas already proved, so each
has degree $-1$ and shifts $(0,-1/2)$.

The component counts of $G^0,G^1,G^2,G^3$ are respectively
$2,1,2,3$.  At $G^3$ the rank-$0$ connections form the main component;
the fixed ranks $2,3$ of $q_2$ give the other two components.  The
selected O-switch graph has one edge from the main component to each
of them, and its reduced incidence matrix is $I_2$.  The determinant
pairing therefore proves nonvanishing for the two prescribed
contractions.  The marked cancellation algorithm removes both copies
of $p$ from each factor.  It leaves the active words $\beta_1$ and
$\sigma_0$, together with exactly the two retained ordinary
$2\times2$ auxiliary blocks.  Applying the proved four-mark collapse
and singleton contraction to each block, and then the simple-factor
functional after the seam--tensor comparison, gives the stated local
target.  This uses the map construction of Statement E, with its
explicit degree $-1$ and $+1$ intermediate maps.

The count is
\[
 r=2,\qquad s=(S-1)(r-1)=3,\qquad
 U=(S-a_1)+B(S-2)=2,\qquad s-U=1=\mu(V).
\]
There is one X switch of degree $0$ and second shift $1/2$, two O
switches of degree $-1$ and second shift $-1/2$, and two prescribed
auxiliary contractions of degree $+1$ and second shift $1/2$.
The four-mark collapses each have degree and shift zero.  Expanding
their formulas adds two further O switches and two further singleton
contractions; thus the completely expanded map uses four of each,
still with one X switch.  These added pairs cancel in both degree
and shift.  Consequently
the complete map has degree $0$, shifts $(0,1/2)$, and sends the
canonical class to the canonical class.

Finally choose the local prepared grid $D_1$ explicitly.  In the
outgoing-column tie order of $\operatorname{wordGrid}(\beta_1)$,
its closing rank-$0$ arc is in row $b=4$, ending in column $0$.
X-stabilize there by inserting a column immediately before column
$0$ and a row immediately after row $4$, with the pattern
$O(W,t),X(W,4),X(0,t)$.  Cut at old column $0$ and the new row $t$.
This produces the following size-$46$ tuple.  Its first X and last O
are in row $0$, its last X is in row $45$, and reversing that single
stabilization proves $D_1\sim\operatorname{wordGrid}(\beta_1)$.
Since $N=1$, the target iterated seam is exactly this grid.

\paragraph{Target prepared grid $D_1$: size $46$.}\mbox{}
\[
\begin{array}{c|rrrrrrrrrrrr}
c&0&1&2&3&4&5&6&7&8&9&10&11\\
O(c)&1&41&6&16&7&21&31&22&36&2&42&8\\
X(c)&0&15&1&6&30&16&21&40&31&41&7&2\\
\hline
c&12&13&14&15&16&17&18&19&20&21&22&23\\
O(c)&17&9&23&32&24&37&3&43&10&18&11&25\\
X(c)&8&22&17&23&36&32&42&9&3&10&24&18\\
\hline
c&24&25&26&27&28&29&30&31&32&33&34&35\\
O(c)&33&26&38&4&44&12&19&13&27&34&28&39\\
X(c)&25&37&33&43&11&4&12&26&19&27&38&34\\
\hline
c&36&37&38&39&40&41&42&43&44&45&&\\
O(c)&5&45&14&20&15&29&35&30&40&0&&\\
X(c)&44&13&5&14&28&20&29&39&35&45&&\\
\end{array}
\]

All tables and intermediate permutations above are generated by the
accompanying elementary script.  It checks every permutation and
O/X inequality, computes the component cycles as
$X^{-1}\circ O$, verifies the displayed row operations and the
strict noninterleaving inequalities, and checks the target frame
coordinates.  The chain-map, degree, filtration, and class claims
follow from the proved general formulas, rather than from a numerical
test of the $120!$-state complex.

\section{Exceptional cases and exact scope}
\label{sec:exceptions}

\begin{longtable}{>{\raggedright\arraybackslash}p{.24\textwidth}>{\raggedright\arraybackslash}p{.69\textwidth}}
\toprule
Case from the brief & Treatment in the construction\\
\midrule
\endhead
One central branch, $b=1$ &
Nothing in the scissor or class argument requires multiple boundary
components.  The formula gives $\mu(V)=2\delta(V)$, so the X count
is even as required.\\[5pt]
One singular point, $N=1$ &
There can still be $B>0$ simple factors.  The construction first
handles all $r=1+B$ factors and then removes the simple tensor
factors.  The final iterated seam is just $D_1$.  The distinct case
$r=1$ is treated explicitly in the proof.\\[5pt]
Several tangent directions &
Choose a central transverse direction first, construct the finite
representative and the global section, then choose a sufficiently
small allowed level.  A small projection shear avoids the finite
set of forbidden central and local tangent directions without
changing that level.  Higher ordinary ramification and repeated
critical values are handled explicitly; simplicity is not assumed.\\[5pt]
Empty $\mathbf w$, ordinary $S$-fold central point &
The word $\mathbf d_S$ still contains every generator, so every
combinatorial rule applies.  Under the full assembly identities
the further condition $b=S$ gives
$\mu(V)=2\delta(V)-S+1=1-M<0$, since $M\ge2$; hence that entire
ordinary-point specialization cannot occur for such an assembly.
This elementary numerical observation is not an invocation of
the conjecture.\\[5pt]
Smallest $S=4$ &
No estimate is weakened at this value.  In the $\mu(V)=1$ case
there is exactly one ordinary X identity, with the O/contraction
pairs having total shift zero.  Section~\ref{sec:worked-four}
gives all relevant indexed tuples for a finite example.\\[5pt]
Missing generators &
The original source contains every generator because of its full
twist.  The external guard retains every generator throughout
initial signed-word normalization.  Later unused ranks are kept
as explicit two-row, two-column components of the defined signed
extension and are removed only by the proved auxiliary maps.\\[5pt]
Grid sizes and successor sizes &
The positive word grids have size $3\ell\ge3$.  Every prepared
factor has one additional row and column.  A cyclic seam
destabilization is used only with at least two factors, so its
temporary size is at least two.  Auxiliary contractions retain the
nontrivial main factor.  No size-zero or size-one complex occurs.
Every final factor size is a successor as required.\\[5pt]
$U=0$, $B=0$, or $s=0$ &
Products over empty sets are identities and the empty incidence
determinant is one.  No auxiliary map is used when $U=0$, and no
functional is used when $B=0$.  The case $s=0$ is exactly the
separately handled $r=1$ case.\\[5pt]
Large $\sum_j a_j-(N-1)$ &
The construction never assumes this quantity is at most $S$.
Each individual $a_j\le S$ follows from the finite projection.
Any excess in the sum is absorbed by the actual count
$\sum_j a_j+B-N-S+1=\mu(V)$.\\
\bottomrule
\end{longtable}

\section{Dependency ledger and references}
\label{sec:ledger}

All entries labelled new below are proved in this document.
The names are proposed short names for later formalization; no claim
of a completed Lean formalization is made.

\begin{longtable}{>{\raggedright\arraybackslash}p{.32\textwidth}>{\raggedright\arraybackslash}p{.45\textwidth}>{\raggedright\arraybackslash}p{.15\textwidth}}
\toprule
Result & Proposed short name and type & Status\\
\midrule
\endhead
Nonvanishing implies the class identity &
\texttt{polynomialImageNonzero}; uses canonical image & Proved\\[3pt]
Composition along ordinary moves &
\texttt{movesToNativeMap}; uses primitive moves and corner reduction & Proved\\[3pt]
Finite projection and event factorization &
\texttt{finiteProjectionEventFactorization}; new analytic & Proved\\[3pt]
O derivatives, saddle pairing, positive-degree vanishing, incidence formula &
\texttt{nativeOSaddleIncidence}; new combinatorial & Proved\\[3pt]
Explicit incoming and outgoing polynomial projection formulas &
\texttt{nativeTopDomainMaps}; new combinatorial & Proved\\[3pt]
Colour naturality of native moves &
\texttt{nativeMoveColourNaturality}; new combinatorial & Proved\\[3pt]
Singleton contraction, transport, and substitution &
\texttt{nativeAuxiliaryContraction}; new combinatorial & Proved\\[3pt]
Paired O switch and contraction &
\texttt{pairedAuxiliaryContraction}; new combinatorial & Proved\\[3pt]
Exact transport and deletion of a singleton &
\texttt{auxiliarySingletonTransport}; new combinatorial & Proved\\[3pt]
Four-mark block normalization &
\texttt{fourMarkBlockCollapse}; new combinatorial & Proved\\[3pt]
Signed relations, shared endpoints, subdivisions, and guard &
\texttt{signedWordGridRelations}; new combinatorial & Proved\\[3pt]
Marked inverse cancellation with a protected grid &
\texttt{markedInverseCancellation}; new combinatorial & Proved\\[3pt]
Common strand by finite Hurwitz moves &
\texttt{commonStrandHurwitz}; new combinatorial & Proved\\[3pt]
Star scissors and explicit rooted forest &
\texttt{scissorIncidenceForest}; new combinatorial & Proved\\[3pt]
Triple preparation, strict seam coordinates, and cyclic exposure &
\texttt{cyclicSeamCoordinates}; new combinatorial & Proved\\[3pt]
Successive factor cleanup and exact auxiliary deletions &
\texttt{seamFactorCleanup}; new combinatorial & Proved\\[3pt]
Simple-factor functional and binary seam reduction &
\texttt{simpleFactorProjection}; new combinatorial, uses binary seam & Proved\\[3pt]
Exact ramification count &
\texttt{exteriorRamificationCount}; uses the word lengths and invariant identities & Proved\\[3pt]
Ordered interchange of common-strand factors &
\texttt{orderedCommonStrandHurwitz}; new combinatorial & Proved\\[3pt]
Statement E &
\texttt{ExteriorNativeMapStatement}; uses the preceding results
and the admitted inputs & Proved\\
\bottomrule
\end{longtable}

The external source ledger is short.  The attached brief~\cite{brief}
is used for exactly the inputs listed in Section~\ref{sec:statement}:
assembly and invariant facts; the central and branch event words;
the native complex and canonical-image theorem; the supplied
commutation and corner-reduction facts; and the binary seam--tensor
equivalence.  These are hypotheses supplied for the problem.
The formulas for incoming and outgoing local domains were checked
against Definitions~13.3.7, 13.3.10, 13.4.2, and~13.4.3 of
Ozsv\'ath--Stipsicz--Szab\'o~\cite{OSS}; their chain, grading, colour,
and class properties needed here are verified directly in this
document.  No theorem identifying grids of isotopic links, no
Markov or Reidemeister theorem, no surface movie theorem, and no
comparison with holomorphic knot Floer homology enters the proof.

\begin{thebibliography}{9}
\bibitem{brief}
\emph{Problem brief: the exterior native map in the Lean proof of
L\^e's conjecture}, supplied Markdown attachment, prepared
2 November 2026.  The input assumptions and the target Statement E
are those of Sections~2, 3.1, and~5.
\bibitem{OSS}
P.~S. Ozsv\'ath, A.~I. Stipsicz, and Z.~Szab\'o,
\emph{Grid Homology for Knots and Links},
American Mathematical Society, 2015.
Author-hosted preliminary text:
\url{https://web.math.princeton.edu/~petero/GridHomologyBook.pdf}.
Only the explicitly identified local domain definitions are used
for a convention check; the needed identities are proved here.
\end{thebibliography}


\end{document}
